% Kaappi: A Scheme Programming Language
% Compiled with XeLaTeX

\documentclass[11pt,twoside,openright]{memoir}

% Load preamble files
% colors.tex - Color definitions for Kaappi book

\usepackage{xcolor}

% Primary accent
\definecolor{darkbrown}{RGB}{64, 64, 64}

% Note box
\definecolor{notebg}{RGB}{245, 245, 245}
\definecolor{noteborder}{RGB}{140, 140, 140}

% Code box
\definecolor{codebg}{RGB}{248, 248, 248}
\definecolor{codeborder}{RGB}{160, 160, 160}

% Warning box
\definecolor{warningbg}{RGB}{255, 255, 255}
\definecolor{warningborder}{RGB}{80, 80, 80}

% Exercise box
\definecolor{exercisebg}{RGB}{238, 238, 238}
\definecolor{exerciseborder}{RGB}{120, 120, 120}

% Output/REPL box
\definecolor{outputbg}{RGB}{235, 235, 235}
\definecolor{outputborder}{RGB}{110, 110, 110}

% Syntax highlighting
\definecolor{keywordcolor}{RGB}{50, 50, 50}
\definecolor{stringcolor}{RGB}{90, 90, 90}
\definecolor{commentcolor}{RGB}{130, 130, 130}

% packages.tex - Package loading for Kaappi book

% Typography
\usepackage[
    protrusion=true,
    expansion=false,
    tracking=false,
    final
]{microtype}
\usepackage{enumitem}
\usepackage{fancyvrb}

% Graphics
\usepackage{tikz}
\usetikzlibrary{positioning, shapes.multipart, arrows.meta}
\usepackage{graphicx}
\usepackage{float}

% Boxes and frames
\usepackage[framemethod=tikz]{mdframed}
\usepackage{adjustbox}

% Code listings
\usepackage{listings}

% Cross-references and navigation
\usepackage{hyperref}
\hypersetup{
    colorlinks=true,
    linkcolor=darkbrown,
    urlcolor=darkbrown,
    bookmarks=true,
    bookmarksnumbered=true,
    pdfauthor={Baiju Muthukadan},
    pdftitle={Kaappi: A Scheme Programming Language},
    pdfsubject={Scheme programming with Kaappi},
}

% Index
\usepackage{imakeidx}
\makeindex[columns=2, title=Index, intoc]

% TOC depth
\setcounter{tocdepth}{2}
\setcounter{secnumdepth}{2}

% fonts.tex - Font configuration using fontspec

\usepackage{fontspec}

\defaultfontfeatures{Ligatures=TeX}

% Main body font: EB Garamond
\setmainfont{EBGaramond}[
    Path = fonts/,
    UprightFont = *-Regular,
    BoldFont = *-Bold,
    ItalicFont = *-Italic,
    BoldItalicFont = *-BoldItalic,
    Extension = .ttf,
    Numbers = Lining,
    Ligatures = {Common, TeX},
]

% Sans-serif for headings: Source Sans 3
\setsansfont{SourceSans3}[
    Path = fonts/,
    UprightFont = *-Regular,
    BoldFont = *-Bold,
    ItalicFont = *-Italic,
    BoldItalicFont = *-BoldItalic,
    Extension = .ttf,
    Numbers = Lining,
    Ligatures = {Common, TeX},
]

% Monospace for code: Fira Code
\setmonofont{FiraCode}[
    Path = fonts/,
    UprightFont = *-Regular,
    BoldFont = *-Bold,
    Extension = .ttf,
    Scale = 0.85,
]

% Math font matched to EB Garamond. Bundled locally (fonts/) for build
% reproducibility rather than relying on the TeX Live copy. OFL 1.1, same
% license as the text fonts. unicode-math must load after the text fonts.
\usepackage{unicode-math}
\setmathfont{Garamond-Math.otf}[
    Path = fonts/,
]

% fonts-memoir.tex - Memoir-specific heading font styles

\renewcommand{\chaptitlefont}{\Huge\sffamily\bfseries\color{darkbrown}}
\renewcommand{\chapnumfont}{\Huge\sffamily\bfseries\color{darkbrown}}
\setsecheadstyle{\Large\sffamily\bfseries\color{darkbrown}}
\setsubsecheadstyle{\large\sffamily\bfseries\color{darkbrown}}
\setsubsubsecheadstyle{\normalsize\sffamily\bfseries\color{darkbrown}}

\renewcommand{\parttitlefont}{\Huge\sffamily\bfseries\color{darkbrown}}
\renewcommand{\partnumfont}{\huge\sffamily\bfseries\color{darkbrown}}

% decorations.tex - Minimal decorative elements for the Kaappi book

\usetikzlibrary{calc}

%% LAMBDA ORNAMENT (section breaks)
\newcommand{\lambdaornament}[1][1]{%
    \begin{tikzpicture}[scale=#1, baseline=-2pt]
        \node[font=\Large, text=darkbrown!70] at (0,0) {$\lambda$};
    \end{tikzpicture}%
}

%% SECTION BREAK
\newcommand{\sectionbreak}{%
    \par\vspace{1.5\baselineskip}
    \noindent\hfill
    \rule{1cm}{0.4pt}\hspace{0.8em}%
    \lambdaornament[0.8]%
    \hspace{0.8em}\rule{1cm}{0.4pt}%
    \hfill\null
    \par\vspace{1.5\baselineskip}
}

%% CHAPTER END ORNAMENT
\newcommand{\chapterend}{%
    \par\vspace{2\baselineskip}
    \noindent\hfill
    \rule{2cm}{0.4pt}\hspace{1em}%
    \lambdaornament%
    \hspace{1em}\rule{2cm}{0.4pt}%
    \hfill\null
    \par\vspace{\baselineskip}
}

% layout.tex - 6" x 9" page layout

% Page dimensions for 6" x 9" (152.4mm x 228.6mm)
\setstocksize{228.6mm}{152.4mm}
\settrimmedsize{228.6mm}{152.4mm}{*}

% Margins: inner 20mm (0.79"), outer 15mm (0.59")
% KDP 307pp minimum: gutter 0.625" (15.875mm), outside/top/bottom 0.25" (6.35mm)
\setlrmarginsandblock{20mm}{15mm}{*}
% Top 15mm, bottom 20mm (extra bottom space for footer page numbers)
\setulmarginsandblock{15mm}{20mm}{*}

% Header and footer spacing
\setheadfoot{14pt}{22pt}
\setheaderspaces{*}{10pt}{*}

\checkandfixthelayout

% TOC page number box: default 1.55em is too narrow for bold 3-digit numbers
\makeatletter
\renewcommand{\@pnumwidth}{2em}
\renewcommand{\@tocrmarg}{2.5em}
\makeatother

% Chapter style
\makeatletter
\makechapterstyle{kaappi}{%
    \chapterstyle{default}
    \renewcommand*{\chapterheadstart}{\vspace*{30pt}}
    \renewcommand*{\afterchapternum}{\par\nobreak\vskip 8pt}
    \renewcommand*{\afterchaptertitle}{%
        \par\nobreak\vskip 8pt
        \noindent\rule{\linewidth}{0.4pt}%
        \par\nobreak\vskip 25pt
    }
    \renewcommand*{\printchaptername}{}
    \renewcommand*{\chapternamenum}{}
    \renewcommand*{\printchapternum}{%
        \centering
        {\chapnumfont \@chapapp\ \thechapter}%
    }
    \renewcommand*{\printchaptertitle}[1]{%
        \centering\chaptitlefont ##1
    }
}
\makeatother
\chapterstyle{kaappi}

% Part style
\renewcommand{\partnamefont}{\normalfont\huge\sffamily\bfseries\color{darkbrown}}
\renewcommand{\partnumfont}{\normalfont\huge\sffamily\bfseries\color{darkbrown}}
\renewcommand{\parttitlefont}{\normalfont\Huge\sffamily\bfseries\color{darkbrown}}
\renewcommand{\printpartname}{Part}
\renewcommand{\printparttitle}[1]{\parttitlefont #1}

% Headers and footers
\makepagestyle{kaappibook}
\makeevenhead{kaappibook}{%
    \small\sffamily\thepage%
}{}{%
    \small\sffamily\leftmark%
}
\makeoddhead{kaappibook}{%
    \small\sffamily\rightmark%
}{}{%
    \small\sffamily\thepage%
}
\makeevenfoot{kaappibook}{}{}{}
\makeoddfoot{kaappibook}{}{}{}
\makeheadrule{kaappibook}{\textwidth}{0.4pt}
\pagestyle{kaappibook}

% Chapter pages use plain style
\aliaspagestyle{chapter}{plain}

% Paragraph settings
\setlength{\parindent}{1.5em}
\setlength{\parskip}{0pt}

% Widow and orphan control
\widowpenalty=10000
\clubpenalty=10000

\frenchspacing
\flushbottom

% Hyphenation control
\doublehyphendemerits=10000
\finalhyphendemerits=5000
\tolerance=1500
\emergencystretch=1em

% styles.tex - Custom environments for the Kaappi book

%% NOTE BOX
\newmdenv[
    backgroundcolor=notebg,
    linecolor=noteborder,
    linewidth=0.5pt,
    innerleftmargin=12pt,
    innerrightmargin=12pt,
    innertopmargin=10pt,
    innerbottommargin=10pt,
    leftmargin=0pt,
    rightmargin=0pt,
    skipabove=\baselineskip,
    skipbelow=\baselineskip,
    roundcorner=2pt,
    topline=true,
    bottomline=true,
    leftline=false,
    rightline=false,
]{notebox}

\newenvironment{note}[1][Note]{%
    \begin{notebox}
    \noindent\textbf{\sffamily\color{darkbrown}#1.}\space
}{%
    \end{notebox}
}

%% WARNING BOX
\newmdenv[
    backgroundcolor=warningbg,
    linecolor=warningborder,
    linewidth=2.5pt,
    roundcorner=0pt,
    innerleftmargin=12pt,
    innerrightmargin=12pt,
    innertopmargin=10pt,
    innerbottommargin=10pt,
    leftmargin=0pt,
    rightmargin=0pt,
    skipabove=\baselineskip,
    skipbelow=\baselineskip,
    leftline=true,
    rightline=false,
    topline=false,
    bottomline=false,
]{warningbox}

\newenvironment{warning}[1][Warning]{%
    \begin{warningbox}
    \noindent\textbf{\sffamily\color{warningborder}#1.}\space
}{%
    \end{warningbox}
}

%% EXERCISE BOX
\newcounter{exercise}[chapter]
\renewcommand{\theexercise}{\thechapter.\arabic{exercise}}

\newmdenv[
    backgroundcolor=white,
    linecolor=exerciseborder,
    linewidth=2pt,
    roundcorner=0pt,
    innerleftmargin=12pt,
    innerrightmargin=12pt,
    innertopmargin=10pt,
    innerbottommargin=10pt,
    leftmargin=0pt,
    rightmargin=0pt,
    skipabove=\baselineskip,
    skipbelow=\baselineskip,
    leftline=true,
    rightline=false,
    topline=false,
    bottomline=false,
]{exercisebox}

\newenvironment{exercise}[1][]{%
    \refstepcounter{exercise}%
    \begin{exercisebox}
    \noindent\textbf{\sffamily\color{exerciseborder}Exercise \theexercise%
    \ifx&#1&\else: #1\fi}\par\vspace{0.5\baselineskip}
}{%
    \end{exercisebox}
}

%% CODE OUTPUT / REPL BOX
\newmdenv[
    backgroundcolor=outputbg,
    linecolor=outputborder,
    linewidth=0.5pt,
    innerleftmargin=12pt,
    innerrightmargin=12pt,
    innertopmargin=10pt,
    innerbottommargin=8pt,
    leftmargin=0pt,
    rightmargin=0pt,
    skipabove=\baselineskip,
    skipbelow=\baselineskip,
    roundcorner=2pt,
]{outputbox}

\newenvironment{codeoutput}[1][]{%
    \begin{outputbox}
    \ifx&#1&\else\noindent\textbf{\sffamily\small\color{darkbrown}#1}\par\vspace{0.3\baselineskip}\fi
    \ttfamily\small
}{%
    \end{outputbox}
}

% listings.tex - Scheme code listing configuration

\lstdefinelanguage{Scheme}{
    % Policy: bold syntax (special forms and library declarations) only.
    % Ordinary procedures (car, map, display, ...) stay unbolded so the
    % highlighting marks what is evaluated specially, not what is common.
    morekeywords={
        define, lambda, if, cond, case, else, and, or,
        when, unless, begin, set!, quote, quasiquote,
        unquote, unquote-splicing,
        let, let*, letrec, letrec*, let-values, let*-values,
        define-values, do, case-lambda,
        define-syntax, syntax-rules, let-syntax, letrec-syntax,
        define-library, import, export, include,
        define-record-type, guard,
        delay, delay-force, parameterize,
        cond-expand, only, except, rename, prefix,
    },
    sensitive=true,
    morecomment=[l]{;},
    morecomment=[s]{\#|}{|\#},
    morestring=[b]{"},
    alsoother={\#},
    literate=
        {->}{$\rightarrow$}1,
}

\lstdefinestyle{scheme}{
    language=Scheme,
    basicstyle=\ttfamily\small,
    keywordstyle=\bfseries\color{keywordcolor},
    commentstyle=\itshape\color{commentcolor},
    stringstyle=\color{stringcolor},
    backgroundcolor=\color{codebg},
    frame=single,
    framerule=0.4pt,
    rulecolor=\color{codeborder},
    xleftmargin=4pt,
    xrightmargin=4pt,
    framexleftmargin=4pt,
    framexrightmargin=4pt,
    aboveskip=\baselineskip,
    belowskip=\baselineskip,
    columns=fullflexible,
    keepspaces=true,
    showstringspaces=false,
    tabsize=2,
    breaklines=true,
    breakatwhitespace=true,
    numbers=none,
    escapeinside={(*@}{@*)},
}

\lstdefinestyle{repl}{
    style=scheme,
    backgroundcolor=\color{outputbg},
    rulecolor=\color{outputborder},
}

\lstset{style=scheme}

% Shorthand for inline code
\newcommand{\code}[1]{\lstinline[style=scheme]{#1}}


\begin{document}

% Front matter
\frontmatter
\pagestyle{empty}

% Title page
\thispagestyle{empty}
\begin{center}
    \vspace*{3cm}
    {\fontsize{32}{38}\selectfont\sffamily\bfseries\color{darkbrown}
     Kaappi}\\[1cm]
    {\Large\sffamily A Scheme Programming Language}\\[3cm]
    {\large\sffamily Baiju Muthukadan}\\[1cm]
    \vfill
    {\normalsize\sffamily R7RS-small $\bullet$ Zig $\bullet$ Batteries Included}\\[1.5cm]
\end{center}
\newpage

% Copyright page
\thispagestyle{empty}
\vspace*{\fill}
\begin{center}
    \textit{Kaappi: A Scheme Programming Language}\\[1em]
    Baiju Muthukadan\\[2em]
    First Edition, 2026\\[2em]
    Copyright \textcopyright{} 2026 Baiju Muthukadan\\
    All rights reserved.\\[2em]
    \vspace{2em}
    \textbf{Draft Edition}\\[0.5em]
    This book is a work in progress.\\
    Feedback is welcome at \href{mailto:baiju.m.mail@gmail.com}{baiju.m.mail@gmail.com}.
\end{center}
\vspace*{\fill}
\cleardoublepage

% Table of Contents
\pagestyle{plain}
\tableofcontents*
\cleardoublepage

% Preface
\chapter{Preface}

This book teaches you how to program in Scheme using Kaappi, a modern implementation that conforms to the R7RS-small standard. It is written for programmers who already know at least one language---such as Python, JavaScript, or Ruby---and want to explore the elegance of Scheme without drowning in academic jargon.

Kaappi provides a complete, batteries-included environment: a fast bytecode compiler, an interactive REPL, over 600 built-in procedures, 72 SRFIs, a package manager, and even a foreign function interface for calling C libraries. This book will guide you from your first expression all the way to writing macros, using continuations, and building real programs.

\section*{How to Use This Book}

The book is organized into four parts:

\begin{description}[leftmargin=1em, labelindent=0em]
    \item[Part I: Getting Started] covers installation, the REPL, and basic expressions.
    \item[Part II: The Language] teaches the core of Scheme: data types, functions, control flow, and I/O.
    \item[Part III: Advanced Topics] introduces macros, modules, error handling, and continuations.
    \item[Part IV: Beyond R7RS] explores Kaappi-specific features like the FFI, fibers, and the SRFI ecosystem.
\end{description}

Each chapter builds on the previous ones, but experienced programmers can skip ahead to topics of interest. Code examples are meant to be typed into the Kaappi REPL as you read.

\cleardoublepage

% Main matter
\mainmatter
\pagestyle{kaappibook}

% Part I: Getting Started
\part{Getting Started}

\chapter{Introduction}

Scheme is one of the oldest and most influential programming languages in the Lisp family. Designed for clarity and minimalism, it has shaped how we think about computation, recursion, and language design itself.

Kaappi is a modern, complete implementation of R7RS-small Scheme, written in Zig. It provides everything you need to write real programs: a fast bytecode compiler, an interactive REPL with syntax highlighting and tab completion, over 600 built-in procedures, 72 SRFIs, a C foreign function interface, green threads, and a package manager.

This chapter explains what Scheme is, why it matters, and what makes Kaappi a good choice for learning and using it.

\chapter{Installation and Setup}

Before writing any Scheme code, you need to install Kaappi on your system. This chapter covers installation on macOS and Linux, verifying your setup, configuring your editor, and getting comfortable with the REPL.

\chapter{First Steps}
\index{expressions}
\label{ch:first-steps}

With Kaappi installed and the REPL ready, it is time to write your first Scheme code. This chapter covers the fundamental building blocks: expressions, basic arithmetic, variables, and defining simple functions. Open a REPL and type along---every example here is meant to be tried interactively.

\section{Everything Is an Expression}

In Python, you distinguish between \emph{statements} (which do things) and \emph{expressions} (which produce values). In Scheme, there is only one category: everything is an expression, and every expression produces a value.

The syntax is uniform. An expression is either an atom (a number, string, or symbol standing alone) or a list enclosed in parentheses\index{S-expression}. When the evaluator sees a list, it treats the first element as the operator and the rest as operands:

\begin{lstlisting}[style=repl]
> (+ 2 3)
5
> (* 6 7)
42
> (- 100 1)
99
\end{lstlisting}

The operator always comes first. This is called \emph{prefix notation}\index{prefix notation}. There is no infix \texttt{2 + 3} in Scheme---it is always \texttt{(+ 2 3)}.

\subsection{No Precedence Rules}

In Python, \texttt{2 + 3 * 4} evaluates to \texttt{14} because \texttt{*} binds tighter than \texttt{+}. In Scheme, you nest expressions explicitly:

\begin{lstlisting}[style=repl]
> (+ 2 (* 3 4))
14
> (* (+ 1 2) (- 10 4))
18
\end{lstlisting}

The parentheses \emph{are} the precedence. There is nothing to memorize, and there is never ambiguity about evaluation order.

\begin{note}[Coming from Python]
If \texttt{(+ 2 3)} looks strange, think of it as a procedure\index{procedure} call where the procedure name goes inside the parentheses: like \texttt{add(2, 3)} but without the commas and with the procedure name moved inside. After a few examples, the pattern becomes natural.
\end{note}

\subsection{Operators Take Multiple Arguments}

Unlike most languages where \texttt{+} is binary, Scheme's arithmetic operators accept any number of arguments:

\begin{lstlisting}[style=repl]
> (+ 1 2 3 4 5)
15
> (* 2 3 4)
24
> (+)
0
> (*)
1
\end{lstlisting}

With zero arguments, \texttt{+} returns 0 (the additive identity) and \texttt{*} returns 1 (the multiplicative identity).

\section{Atoms}
\index{atoms}

Atoms are the simplest values in Scheme---they cannot be broken into smaller parts.

\subsection{Numbers}
\index{numbers}

Kaappi supports a full numeric tower:

\begin{lstlisting}[style=repl]
> 42
42
> -7
-7
> 3.14
3.14
> 1/3
1/3
> (+ 1/3 1/6)
1/2
\end{lstlisting}

Integers and exact rationals stay exact---there is no floating-point drift when you work with fractions. Kaappi also supports arbitrarily large integers (bignums):

\begin{lstlisting}[style=repl]
> (* 1000000000 1000000000 1000000000)
1000000000000000000000000000
\end{lstlisting}

\begin{note}[Coming from Python]
Python 3 also has arbitrary-precision integers, so this will feel familiar. The difference is that Scheme has exact rationals (\texttt{1/3}) as a built-in type---Python requires the \texttt{fractions} module.
\end{note}

\subsection{Strings}
\index{strings}

Strings are delimited by double quotes:\index{string-length}\index{string-append}

\begin{lstlisting}[style=repl]
> "hello, world"
"hello, world"
> (string-length "hello")
5
> (string-append "foo" "bar")
"foobar"
\end{lstlisting}

Single quotes mean something entirely different in Scheme (they denote quotation, which we will cover shortly), so strings are always double-quoted.

\subsection{Booleans}
\index{booleans}

The boolean values are \texttt{\#t} (true) and \texttt{\#f} (false). There is one critical rule that surprises Python programmers: \textbf{only \texttt{\#f} is false}. Everything else---including \texttt{0}, the empty string, and the empty list---is true:

\begin{lstlisting}[style=repl]
> (if 0 "truthy" "falsy")
"truthy"
> (if "" "truthy" "falsy")
"truthy"
> (if '() "truthy" "falsy")
"truthy"
> (if #f "truthy" "falsy")
"falsy"
\end{lstlisting}

\begin{warning}[Only \texttt{\#f} is false]
This is one of the most common sources of bugs when coming from Python. In Scheme, \texttt{(if 0 ...)} takes the ``then'' branch because \texttt{0} is not \texttt{\#f}. If you want to test whether a number is zero, use \texttt{(zero? n)} or \texttt{(= n 0)} explicitly.
\end{warning}

\subsection{Characters}
\index{characters}

Individual characters are written with a \texttt{\#\textbackslash} prefix:

\begin{lstlisting}[style=repl]
> #\a
#\a
> #\space
#\space
> #\newline
#\newline
> (char-alphabetic? #\Z)
#t
\end{lstlisting}

Characters are distinct from single-character strings.

\subsection{Symbols}
\index{symbols}

A symbol is a lightweight, interned name. Symbols are compared by identity (fast pointer comparison), not by character content like strings:

\begin{lstlisting}[style=repl]
> 'hello
hello
> (symbol? 'hello)
#t
> (eq? 'abc 'abc)
#t
\end{lstlisting}

The \texttt{'} (single quote) is called \emph{quote}\index{quote}. It prevents the expression that follows from being evaluated. Without it, \texttt{hello} would be looked up as a variable name, causing an ``unbound variable'' error.

\section{Defining Variables}
\index{define}

Use \texttt{define} to bind a name to a value:

\begin{lstlisting}[style=repl]
> (define pi 3.14159)
> pi
3.14159
> (* pi 10 10)
314.159
\end{lstlisting}

Once defined, you can use the name anywhere an expression is expected. Names in Scheme can contain characters that would be illegal in most languages: hyphens, question marks, and exclamation marks are all common:

\begin{lstlisting}[style=repl]
> (define my-favorite-number 42)
> my-favorite-number
42
\end{lstlisting}

\begin{note}[Naming conventions]
Scheme uses \texttt{kebab-case}\index{kebab-case} (words separated by hyphens) instead of \texttt{snake\_case} or \texttt{camelCase}. Predicates end with \texttt{?} (e.g., \texttt{zero?}, \texttt{string?}). Mutating procedures end with \texttt{!} (e.g., \texttt{set!}, \texttt{vector-set!}).
\end{note}

\section{Defining Functions}
\index{functions}\index{lambda}\index{procedures}

Scheme officially calls these \emph{procedures}\index{procedure} rather than functions, because they can have side effects (like printing output). This book uses both terms interchangeably.

To define a procedure, put the name and parameters in a list after \texttt{define}:

\begin{lstlisting}[style=repl]
> (define (square x) (* x x))
> (square 7)
49
> (square 1.5)
2.25
\end{lstlisting}

This is shorthand for the more explicit form using \texttt{lambda}, which creates an anonymous function:

\begin{lstlisting}[style=repl]
> (define square (lambda (x) (* x x)))
> (square 5)
25
\end{lstlisting}

\texttt{lambda}\index{lambda} is like Python's \texttt{lambda} keyword, but without the one-expression limitation---a Scheme lambda can contain any number of expressions in its body.

Functions can take multiple parameters:

\begin{lstlisting}[style=repl]
> (define (area width height) (* width height))
> (area 5 3)
15
\end{lstlisting}

And they can call other functions, including themselves:

\begin{lstlisting}[style=repl]
> (define (sum-of-squares a b)
    (+ (square a) (square b)))
> (sum-of-squares 3 4)
25
\end{lstlisting}

\subsection{Functions Are Values}

Functions are ordinary values in Scheme. You can pass them as arguments to other functions:

\begin{lstlisting}[style=repl]
> (define (apply-twice f x) (f (f x)))
> (apply-twice square 3)
81
> (apply-twice square 2)
16
\end{lstlisting}

In the first call, \texttt{(square 3)} gives \texttt{9}, then \texttt{(square 9)} gives \texttt{81}. In the second, \texttt{(square 2)} gives \texttt{4}, then \texttt{(square 4)} gives \texttt{16}. This ability to pass functions around is what makes Scheme a \emph{functional} language, and it is the foundation for many powerful patterns you will learn in later chapters.

\section{Basic Conditionals}
\index{if}\index{conditionals}

The \texttt{if} expression has three parts: a test, a ``then'' expression, and an ``else'' expression:

\begin{lstlisting}[style=repl]
> (if (> 5 3) "yes" "no")
"yes"
> (if (< 5 3) "yes" "no")
"no"
\end{lstlisting}

Because \texttt{if} is an expression (not a statement), it produces a value. You can use it anywhere a value is expected:

\begin{lstlisting}[style=repl]
> (define (abs x)
    (if (< x 0) (- x) x))
> (abs -7)
7
> (abs 3)
3
\end{lstlisting}

\begin{note}[Coming from Python]
This is like Python's ternary expression \texttt{"yes" if 5 > 3 else "no"}, but in Scheme the condition comes first. The pattern is always \texttt{(if test then else)}.
\end{note}

For multiple branches, use \texttt{cond}\index{cond}. It works like Python's \texttt{if/elif/else}:

\begin{lstlisting}[style=repl]
> (define (classify n)
    (cond
      ((< n 0) "negative")
      ((= n 0) "zero")
      (else "positive")))
> (classify -5)
"negative"
> (classify 0)
"zero"
> (classify 42)
"positive"
\end{lstlisting}

Each clause is \texttt{(test result)}. The \texttt{else} clause matches when none of the earlier tests are true.

\section{Comments}
\index{comments}

Line comments start with a semicolon and extend to the end of the line:

\begin{lstlisting}
; This is a comment
(define x 10) ; inline comment
\end{lstlisting}

Block comments use \texttt{\#|} and \texttt{|\#}:

\begin{lstlisting}
#|
This is a
multi-line comment
|#
\end{lstlisting}

A datum comment\index{datum comment} \texttt{\#;} comments out the next expression:

\begin{lstlisting}[style=repl]
> (+ 1 #;(this is ignored) 2)
3
\end{lstlisting}

This last form is particularly useful for temporarily disabling an argument or clause without deleting it.

\section{Putting It Together}

Let's write a slightly larger example that uses what we have learned. Here is a function that converts a temperature from Fahrenheit to Celsius:

\begin{lstlisting}[style=repl]
> (define (fahrenheit->celsius f)
    (* (/ 5 9) (- f 32)))
> (fahrenheit->celsius 212)
100
> (fahrenheit->celsius 32)
0
> (fahrenheit->celsius 72)
200/9
\end{lstlisting}

Notice that \texttt{(fahrenheit->celsius 72)} returns the exact rational \texttt{200/9} rather than a floating-point approximation. Kaappi preserves exactness when all inputs are exact. If you want a decimal, use a floating-point literal:

\begin{lstlisting}[style=repl]
> (define (fahrenheit->celsius-approx f)
    (* (/ 5.0 9.0) (- f 32)))
> (fahrenheit->celsius-approx 72)
22.22222222222222
\end{lstlisting}

\section*{Summary}

\begin{itemize}
    \item Scheme programs are built entirely from expressions in prefix notation, so there are no operator-precedence rules to memorize.
    \item The atomic data types are numbers, strings, booleans, characters, and symbols.
    \item \texttt{define} binds both variables and named functions, and functions are themselves first-class values.
    \item \texttt{if} and \texttt{cond} select between alternatives based on a condition.
    \item Code can be documented with line comments (\texttt{;}) and block comments.
\end{itemize}

\section*{Exercises}

\begin{exercise}[Predict and Check]
Before typing each expression into the REPL, predict what the result will be. Then check your answer:

\begin{lstlisting}[style=repl]
> (+ (* 3 5) (- 10 6))
> (/ 1 3)
> (+ 1/4 1/4 1/4 1/4)
> (* (+ 1 2) (+ 3 4) (+ 5 6))
> (- (* 2 3) (* 4 5))
\end{lstlisting}

Which results surprised you? Pay attention to how Kaappi handles exact rationals versus integers.
\end{exercise}

\begin{exercise}[Celsius to Fahrenheit]
The chapter showed a \code{fahrenheit->celsius} function. Write the inverse: a function \code{celsius->fahrenheit} that takes a temperature in Celsius and returns Fahrenheit. The formula is $F = C \times \frac{9}{5} + 32$. Test it in the REPL:

\begin{lstlisting}[style=repl]
> (celsius->fahrenheit 0)
32
> (celsius->fahrenheit 100)
212
> (celsius->fahrenheit 37)
493/5
\end{lstlisting}

In Python, this would be \code{def celsius\_to\_fahrenheit(c): return c * 9/5 + 32}.
\end{exercise}

\begin{exercise}[Absolute Value]
Write a function \code{(absolute-value n)} that returns the absolute value of \code{n} using \code{if}. It should handle positive numbers, negative numbers, and zero:

\begin{lstlisting}[style=repl]
> (absolute-value -7)
7
> (absolute-value 5)
5
> (absolute-value 0)
0
\end{lstlisting}

Compare your solution to the \code{abs} function defined earlier in this chapter. Are they equivalent?
\end{exercise}

\begin{exercise}[Recursive Power]
Write a function \code{(power base exponent)} that computes $\mathit{base}^{\mathit{exponent}}$ using recursion. Assume \code{exponent} is a non-negative integer. The base case is: anything raised to the power 0 is 1. The recursive case is: $b^n = b \times b^{n-1}$.

\begin{lstlisting}[style=repl]
> (power 2 10)
1024
> (power 3 0)
1
> (power 5 3)
125
\end{lstlisting}

In Python, you would write:

\begin{lstlisting}[language=Python, style=scheme, keywordstyle=\bfseries]
def power(base, exponent):
    if exponent == 0:
        return 1
    return base * power(base, exponent - 1)
\end{lstlisting}

Translate this logic into Scheme using \code{if}, \code{=}, and \code{-}.
\end{exercise}

\begin{exercise}[Temperature Classifier]
Write a function \code{(describe-temp t)} that takes a temperature in Celsius and returns a string describing it. Use \code{cond} with these ranges:

\begin{itemize}
    \item Below 0: \code{"freezing"}
    \item 0 to 15: \code{"cold"}
    \item 16 to 25: \code{"warm"}
    \item Above 25: \code{"hot"}
\end{itemize}

\begin{lstlisting}[style=repl]
> (describe-temp -5)
"freezing"
> (describe-temp 10)
"cold"
> (describe-temp 22)
"warm"
> (describe-temp 35)
"hot"
\end{lstlisting}

Hint: you will need both \code{<} and \code{<=} (or \code{>=}) to express the boundary conditions.
\end{exercise}

In the next chapter, we will explore Scheme's data types in depth.

\chapterend


% Part II: The Language
\part{The Language}

\chapter{Data Types}
\index{data types}

Chapter~3 introduced the basic atoms in passing. This chapter examines Scheme's data types in depth: the full numeric tower, strings, characters, booleans, symbols, vectors, and bytevectors. You will learn how to create values of each type, the key procedures for working with them, and how to test what type a value is.

\section{The Numeric Tower}
\index{numbers}\index{numeric tower}

Kaappi supports five kinds of numbers, organized in a hierarchy called the \emph{numeric tower}:

\begin{enumerate}
    \item \textbf{Integers} --- whole numbers, either fixnums (fast, 63-bit) or bignums (arbitrary precision)
    \item \textbf{Rationals} --- exact fractions like \texttt{1/3} or \texttt{-7/4}
    \item \textbf{Real numbers} --- flonums (IEEE 754 double-precision floating-point)
    \item \textbf{Complex numbers} --- pairs of reals, like \texttt{3+4i}
\end{enumerate}

Every integer is also a rational (it is \texttt{n/1}). Every rational is also a real. Every real is also a complex (with zero imaginary part). The arithmetic procedures work uniformly across all levels.

\subsection{Integers}
\index{integers}\index{fixnum}\index{bignum}

Small integers (up to 63 bits) are stored as \emph{fixnums}---unboxed values that require no heap allocation. When a fixnum overflows, Kaappi automatically promotes it to a \emph{bignum} with unlimited precision:

\begin{lstlisting}[style=repl]
> 42
42
> -100
-100
> (expt 2 100)
1267650600228229401496703205376
> (* 999999999999 999999999999)
999999999998000000000001
\end{lstlisting}

There is no silent wraparound. If the result does not fit in a fixnum, it becomes a bignum transparently.

\subsection{Exact Rationals}
\index{rationals}

When you divide two integers and the result is not an integer, Kaappi returns an exact rational:

\begin{lstlisting}[style=repl]
> (/ 1 3)
1/3
> (/ 22 7)
22/7
> (+ 1/3 1/6)
1/2
> (* 3 1/3)
1
\end{lstlisting}

Rationals are always reduced to lowest terms. They stay exact through arbitrarily long chains of computation---no floating-point drift.

\begin{note}[Coming from Python]
Python's \texttt{/} always returns a float. To get exact fractions in Python, you need \texttt{from fractions import Fraction}. In Scheme, exact arithmetic is the default.
\end{note}

\subsection{Floating-Point Numbers}
\index{floating-point}\index{flonum}

Floating-point numbers (flonums) are written with a decimal point or in scientific notation:

\begin{lstlisting}[style=repl]
> 3.14
3.14
> 2.5e10
25000000000.0
> (/ 1.0 3.0)
0.3333333333333333
\end{lstlisting}

\subsection{Exact vs.\ Inexact}
\index{exact}\index{inexact}

Scheme distinguishes between \emph{exact} and \emph{inexact} numbers. Integers and rationals are exact; flonums are inexact. Mixing them produces an inexact result:

\begin{lstlisting}[style=repl]
> (+ 1/3 0.1)
0.43333333333333335
> (exact? 1/3)
#t
> (inexact? 3.14)
#t
> (exact? 42)
#t
\end{lstlisting}

You can convert between exact and inexact representations:

\begin{lstlisting}[style=repl]
> (inexact 1/3)
0.3333333333333333
> (exact 0.5)
1/2
\end{lstlisting}

\subsection{Complex Numbers}
\index{complex numbers}

Complex numbers are written in rectangular form with \texttt{+} or \texttt{-} separating the real and imaginary parts:

\begin{lstlisting}[style=repl]
> 3+4i
3+4i
> (+ 1+2i 3+4i)
4+6i
> (magnitude 3+4i)
5
> (sqrt -1)
0+1i
\end{lstlisting}

\subsection{Numeric Predicates and Comparisons}

Scheme provides predicates to classify numbers and comparison procedures that work across the numeric tower:

\begin{lstlisting}[style=repl]
> (number? 42)
#t
> (integer? 3.0)
#t
> (rational? 1/3)
#t
> (zero? 0)
#t
> (positive? 5)
#t
> (negative? -3)
#t
> (= 1 1.0)
#t
> (< 1 2 3 4 5)
#t
\end{lstlisting}

Note that \texttt{<}, \texttt{>}, \texttt{<=}, and \texttt{>=} accept multiple arguments and check that the entire sequence is ordered.

\subsection{Common Arithmetic Procedures}

\begin{lstlisting}[style=repl]
> (+ 1 2 3)
6
> (- 10 3)
7
> (* 2 3 4)
24
> (/ 10 3)
10/3
> (quotient 10 3)
3
> (remainder 10 3)
1
> (modulo 10 3)
1
> (abs -7)
7
> (max 3 1 4 1 5)
5
> (min 3 1 4 1 5)
1
> (expt 2 10)
1024
> (sqrt 16)
4
\end{lstlisting}

\section{Strings}
\index{strings}

Strings are sequences of Unicode characters, enclosed in double quotes. Kaappi stores strings as UTF-8 internally but indexes them by codepoint position:

\begin{lstlisting}[style=repl]
> "hello, world"
"hello, world"
> (string-length "hello")
5
> (string-ref "hello" 0)
#\h
> (string-ref "hello" 4)
#\o
\end{lstlisting}

\subsection{String Operations}

\begin{lstlisting}[style=repl]
> (string-append "foo" "bar" "baz")
"foobarbaz"
> (substring "hello, world" 7 12)
"world"
> (string-upcase "hello")
"HELLO"
> (string-downcase "HELLO")
"hello"
> (string-contains "hello world" "world")
6
\end{lstlisting}

\subsection{String Conversion}

\begin{lstlisting}[style=repl]
> (number->string 42)
"42"
> (number->string 3.14)
"3.14"
> (string->number "123")
123
> (string->number "not-a-number")
#f
> (symbol->string 'hello)
"hello"
> (string->symbol "world")
world
\end{lstlisting}

\texttt{string->number} returns \texttt{\#f} (not an error) if the string cannot be parsed as a number. This makes it safe to use on user input.

\subsection{Strings Are Immutable}

In R7RS Scheme, strings are immutable by default. You cannot modify a character in place. Instead, you create new strings from existing ones:

\begin{lstlisting}[style=repl]
> (define greeting "hello")
> (string-append (string-upcase (substring greeting 0 1))
                 (substring greeting 1 5))
"Hello"
\end{lstlisting}

\begin{note}[Coming from Python]
Python strings are also immutable---operations like \texttt{upper()} return new strings. Scheme works the same way.
\end{note}

\section{Characters}
\index{characters}

Characters represent individual Unicode codepoints. They are written with a \texttt{\#\textbackslash} prefix:

\begin{lstlisting}[style=repl]
> #\a
#\a
> #\Z
#\Z
> #\space
#\space
> #\newline
#\newline
> #\tab
#\tab
\end{lstlisting}

Characters are distinct from strings. A character is a single value; a string is a sequence of characters.

\subsection{Character Predicates}

\begin{lstlisting}[style=repl]
> (char-alphabetic? #\a)
#t
> (char-numeric? #\5)
#t
> (char-whitespace? #\space)
#t
> (char-upper-case? #\A)
#t
> (char-lower-case? #\a)
#t
\end{lstlisting}

\subsection{Character Conversion}

\begin{lstlisting}[style=repl]
> (char-upcase #\a)
#\A
> (char-downcase #\Z)
#\z
> (char->integer #\A)
65
> (integer->char 955)
#\x3BB
\end{lstlisting}

The last example produces the Greek lambda character---Kaappi has full Unicode support.

\section{Booleans}
\index{booleans}

Booleans were introduced in Chapter~3, but here is a summary of the key procedures:

\begin{lstlisting}[style=repl]
> (boolean? #t)
#t
> (boolean? 0)
#f
> (not #f)
#t
> (not 0)
#f
> (not '())
#f
\end{lstlisting}

Remember: \texttt{not} returns \texttt{\#t} only when given \texttt{\#f}. Every other value---including \texttt{0}, the empty string, and the empty list---is considered true.

\subsection{Boolean Combinators}

\texttt{and} and \texttt{or} are short-circuiting and return the determining value (not necessarily \texttt{\#t} or \texttt{\#f}):

\begin{lstlisting}[style=repl]
> (and 1 2 3)
3
> (and 1 #f 3)
#f
> (or #f #f 42)
42
> (or #f #f #f)
#f
\end{lstlisting}

\texttt{and} returns the last value if all are true, or the first false value. \texttt{or} returns the first true value, or the last value if all are false.

\begin{note}[Coming from Python]
This is exactly how Python's \texttt{and}/\texttt{or} work: \texttt{1 and 2 and 3} returns \texttt{3}, and \texttt{False or 0 or 42} returns \texttt{42}. Same semantics, different syntax.
\end{note}

\section{Symbols}
\index{symbols}

Symbols are interned identifiers. Two symbols with the same name are always the same object in memory, which makes comparison instantaneous:

\begin{lstlisting}[style=repl]
> (eq? 'hello 'hello)
#t
> (eq? 'hello 'world)
#f
> (symbol? 'foo)
#t
> (symbol->string 'foo)
"foo"
> (string->symbol "bar")
bar
\end{lstlisting}

Symbols are commonly used as keys in association lists (alists), as tags in data structures, and as identifiers in programs that manipulate code.

\begin{lstlisting}[style=repl]
> (define colors '((red . "#ff0000")
                   (green . "#00ff00")
                   (blue . "#0000ff")))
> (assq 'green colors)
(green . "#00ff00")
> (cdr (assq 'green colors))
"#00ff00"
\end{lstlisting}

\section{Vectors}
\index{vectors}

Vectors are fixed-size, indexed collections---like Python lists or arrays. They provide O(1) access by index:

\begin{lstlisting}[style=repl]
> (define v (vector 10 20 30 40 50))
> (vector-ref v 0)
10
> (vector-ref v 4)
50
> (vector-length v)
5
\end{lstlisting}

Vector literals are written with \texttt{\#(}:

\begin{lstlisting}[style=repl]
> #(1 2 3)
#(1 2 3)
\end{lstlisting}

Unlike Scheme lists (which are linked), vectors support efficient random access. Use vectors when you need to look up elements by position.

\subsection{Modifying Vectors}

Vectors are mutable---you can change individual elements:

\begin{lstlisting}[style=repl]
> (define v (vector 1 2 3))
> (vector-set! v 1 99)
> v
#(1 99 3)
\end{lstlisting}

\subsection{Vector Operations}

\begin{lstlisting}[style=repl]
> (vector-map + #(1 2 3) #(10 20 30))
#(11 22 33)
> (vector->list #(a b c))
(a b c)
> (list->vector '(x y z))
#(x y z)
> (make-vector 5 0)
#(0 0 0 0 0)
\end{lstlisting}

\section{Bytevectors}
\index{bytevectors}

Bytevectors are sequences of bytes (unsigned 8-bit integers). They are used for binary data, file I/O, and interfacing with C libraries:

\begin{lstlisting}[style=repl]
> (define bv #u8(72 101 108 108 111))
> (bytevector-length bv)
5
> (bytevector-u8-ref bv 0)
72
> (utf8->string bv)
"Hello"
\end{lstlisting}

You can create bytevectors from strings and vice versa:

\begin{lstlisting}[style=repl]
> (string->utf8 "hello")
#u8(104 101 108 108 111)
> (utf8->string #u8(104 101 108 108 111))
"hello"
\end{lstlisting}

\section{Type Predicates}
\index{predicates}

Scheme provides predicates to test the type of any value. Each returns \texttt{\#t} or \texttt{\#f}:

\begin{lstlisting}[style=repl]
> (number? 42)
#t
> (string? "hello")
#t
> (char? #\a)
#t
> (boolean? #t)
#t
> (symbol? 'foo)
#t
> (vector? #(1 2 3))
#t
> (bytevector? #u8(1 2 3))
#t
> (pair? '(1 2 3))
#t
> (null? '())
#t
> (procedure? +)
#t
\end{lstlisting}

These predicates are essential for writing functions that handle different types of input, and they come up frequently when processing heterogeneous data.

\sectionbreak

You now have a thorough understanding of Scheme's data types. The next chapter covers the most important compound data structure: lists and pairs.

\chapter{Lists and Pairs}

Lists are the fundamental data structure in Scheme. Built from pairs (cons cells), they are used for everything from storing collections to representing code itself. This chapter covers \texttt{cons}, \texttt{car}, \texttt{cdr}, list construction, quoting, and common list operations.

\chapter{Functions}

Functions are the building blocks of Scheme programs. This chapter covers defining functions with \texttt{define} and \texttt{lambda}, writing recursive functions, understanding closures, and how Kaappi guarantees proper tail-call optimization so recursive functions can run in constant stack space.

\chapter{Control Flow}

Scheme's control flow is built on expressions, not statements. This chapter covers conditionals (\texttt{if}, \texttt{cond}, \texttt{case}), boolean combinators (\texttt{and}, \texttt{or}), and convenience forms (\texttt{when}, \texttt{unless}). You will also learn why only \texttt{\#f} is false in Scheme.

\chapter{Local Bindings}
\index{local bindings}
\label{ch:bindings}

In Scheme, local variables are introduced with \emph{binding forms}. Unlike Python, where you create a variable simply by assigning to it, Scheme makes the scope\index{scope} of every variable explicit. This chapter covers \texttt{let}, \texttt{let*}, \texttt{letrec}, internal definitions, and the relationship between \texttt{let} and \texttt{lambda}.

\section{\texttt{let}: Parallel Bindings}
\index{let}

\texttt{let} introduces one or more local variables:

\begin{lstlisting}[style=repl]
> (let ((x 10)
        (y 20))
    (+ x y))
30
\end{lstlisting}

The general form is:

\begin{lstlisting}
(let ((var1 expr1)
      (var2 expr2)
      ...)
  body ...)
\end{lstlisting}

Each \texttt{var} is bound to the value of its \texttt{expr}, and then the \texttt{body} expressions are evaluated in order. The result of the last body expression is the value of the whole \texttt{let}.

\subsection{Bindings Are Parallel}

The key rule: all the \texttt{expr} values are computed \emph{before} any binding takes effect. This means one binding cannot refer to another:

\begin{lstlisting}[style=repl]
> (let ((x 5)
        (y (* x 2)))   ; ERROR: x is not yet bound here
    y)
; Error: unbound variable: x
\end{lstlisting}

This fails because \texttt{x} is not visible to the expression \texttt{(* x 2)}---the bindings are computed in the outer environment. If you need sequential bindings, use \texttt{let*}.

\subsection{Scope Is Limited}

Variables introduced by \texttt{let} exist only within the body. They do not leak into the surrounding scope:

\begin{lstlisting}[style=repl]
> (define x 100)
> (let ((x 5))
    (display x) (newline)
    x)
5
5
> x
100
\end{lstlisting}

The outer \texttt{x} is temporarily shadowed\index{shadowing} inside the \texttt{let}, but remains unchanged afterward.

\subsection{Practical Uses}

Use \texttt{let} to name intermediate results, avoiding repeated computation:

\begin{lstlisting}[style=repl]
> (define (distance x1 y1 x2 y2)
    (let ((dx (- x2 x1))
          (dy (- y2 y1)))
      (sqrt (+ (* dx dx) (* dy dy)))))
> (distance 0 0 3 4)
5
\end{lstlisting}

\section{\texttt{let*}: Sequential Bindings}
\index{let*}

\texttt{let*} is like \texttt{let}, but each binding is visible to subsequent ones. Bindings are evaluated left to right:

\begin{lstlisting}[style=repl]
> (let* ((x 3)
         (y (* x x))
         (z (+ x y)))
    z)
12
\end{lstlisting}

Here, \texttt{y} uses \texttt{x} (which is 3), and \texttt{z} uses both \texttt{x} and \texttt{y}.

\begin{note}[Coming from Python]
\texttt{let*} behaves like consecutive variable assignments in Python:
\begin{Verbatim}[fontsize=\small]
x = 3
y = x * x
z = x + y
\end{Verbatim}
The sequential nature is the same. The difference is that these Scheme bindings are scoped---they vanish after the body, whereas Python variables persist in the enclosing function.
\end{note}

\subsection{When to Use \texttt{let} vs.\ \texttt{let*}}

Use \texttt{let} when bindings are independent of each other. Use \texttt{let*} when later bindings depend on earlier ones. In practice, many programmers default to \texttt{let*} since it is never wrong where \texttt{let} would work, but using \texttt{let} when possible communicates that the bindings are independent.

\section{\texttt{letrec}: Recursive Bindings}
\index{letrec}

\texttt{letrec} allows bindings to refer to each other, enabling mutually recursive local functions:

\begin{lstlisting}[style=repl]
> (letrec ((my-even? (lambda (n)
                       (if (= n 0) #t
                           (my-odd? (- n 1)))))
           (my-odd? (lambda (n)
                      (if (= n 0) #f
                          (my-even? (- n 1))))))
    (list (my-even? 10) (my-odd? 10)))
(#t #f)
\end{lstlisting}

\texttt{my-even?} calls \texttt{my-odd?} and vice versa. Neither \texttt{let} nor \texttt{let*} could express this---\texttt{let} does not allow forward references, and \texttt{let*} only allows references to earlier bindings.

\begin{warning}[Initialization restriction]
The right-hand side expressions in \texttt{letrec} must not evaluate the variables being defined during initialization. In practice, this means the expressions should be \texttt{lambda} forms (which defer evaluation). Do not write \texttt{(letrec ((x (+ y 1)) (y 5)) ...)}---this is undefined behavior.
\end{warning}

\subsection{\texttt{letrec*}}
\index{letrec*}

\texttt{letrec*} is like \texttt{letrec} but guarantees left-to-right evaluation order. It is slightly more permissive: later bindings can use the values of earlier bindings (like \texttt{let*}), and all bindings can be recursive (like \texttt{letrec}):

\begin{lstlisting}[style=repl]
> (letrec* ((x 1)
            (y (+ x 1))
            (f (lambda (n) (+ n x y))))
    (f 10))
13
\end{lstlisting}

\section{Named \texttt{let}: Loops Revisited}
\index{named let}

Named \texttt{let} was introduced in Chapter~\ref{ch:functions}. Here we examine it more carefully as a binding form.

\begin{lstlisting}
(let name ((var init) ...)
  body ...)
\end{lstlisting}

This binds \texttt{name} to a procedure with parameters \texttt{var ...}, calls it with the \texttt{init} values, and makes \texttt{name} available for recursive calls within the body:

\begin{lstlisting}[style=repl]
> (let loop ((numbers '(1 2 3 4 5))
             (sum 0))
    (if (null? numbers)
        sum
        (loop (cdr numbers)
              (+ sum (car numbers)))))
15
\end{lstlisting}

\subsection{Building Lists with Named \texttt{let}}

A common pattern is accumulating results:

\begin{lstlisting}[style=repl]
> (define (squares-up-to n)
    (let loop ((i 1) (acc '()))
      (if (> i n)
          (reverse acc)
          (loop (+ i 1) (cons (* i i) acc)))))
> (squares-up-to 5)
(1 4 9 16 25)
\end{lstlisting}

The accumulator builds the list in reverse (since \texttt{cons} adds to the front), and \texttt{reverse} fixes the order at the end. This is idiomatic Scheme.

\subsection{Named \texttt{let} vs.\ \texttt{do}}

Named \texttt{let} and \texttt{do} can express the same loops. Named \texttt{let} is preferred by most Scheme programmers because the recursive structure is explicit and the control flow is clearer:

\begin{lstlisting}
;; Named let (preferred)
(let loop ((i 0) (sum 0))
  (if (= i 10) sum
      (loop (+ i 1) (+ sum i))))

;; Equivalent do
(do ((i 0 (+ i 1))
     (sum 0 (+ sum i)))
    ((= i 10) sum))
\end{lstlisting}

Both produce 45. Use whichever reads better for your situation.

\section{Internal Definitions}
\index{internal definitions}

Inside a \texttt{lambda}, \texttt{let}, or \texttt{define} body, you can use \texttt{define} to create local bindings:

\begin{lstlisting}[style=repl]
> (define (circle-stats radius)
    (define pi 3.141592653589793)
    (define r-squared (* radius radius))
    (define area (* pi r-squared))
    (define circumference (* 2 pi radius))
    (values area circumference))
> (circle-stats 5)
78.53981633974483
31.41592653589793
\end{lstlisting}

Internal \texttt{define}s are equivalent to \texttt{letrec*}---they can refer to each other and allow mutual recursion. This is the most common way to define local helper functions:

\begin{lstlisting}[style=repl]
> (define (process-list lst)
    (define (helper x)
      (* x x))
    (map helper lst))
> (process-list '(1 2 3 4))
(1 4 9 16)
\end{lstlisting}

\begin{note}[Style preference]
Some programmers prefer internal \texttt{define} for local helpers (reads like top-level code), while others prefer \texttt{let}/\texttt{letrec} (makes scope visually obvious). Both are correct; use what you find clearer.
\end{note}

\section{\texttt{let} Is Syntactic Sugar for \texttt{lambda}}

Understanding the relationship between \texttt{let} and \texttt{lambda} deepens your understanding of both. Every \texttt{let} can be mechanically translated to an immediately-applied \texttt{lambda}:

\begin{lstlisting}
;; These two are identical:
(let ((x 5) (y 10))
  (+ x y))

((lambda (x y) (+ x y)) 5 10)
\end{lstlisting}

The \texttt{let} creates an anonymous function with parameters \texttt{x} and \texttt{y}, then immediately calls it with arguments 5 and 10. This is not just a theoretical curiosity---it explains why \texttt{let} bindings are parallel (arguments to a function call are evaluated independently) and why the body has its own scope.

\section{\texttt{let-values}: Destructuring Multiple Values}
\index{let-values}

When a function returns multiple values (using \texttt{values}), you can bind them with \texttt{let-values}:

\begin{lstlisting}[style=repl]
> (define (min-max lst)
    (let loop ((rest (cdr lst))
               (lo (car lst))
               (hi (car lst)))
      (if (null? rest)
          (values lo hi)
          (loop (cdr rest)
                (min lo (car rest))
                (max hi (car rest))))))
> (let-values (((lo hi) (min-max '(3 1 4 1 5 9 2 6))))
    (display "Min: ") (display lo) (newline)
    (display "Max: ") (display hi) (newline))
Min: 1
Max: 9
\end{lstlisting}

\section{Choosing the Right Binding Form}

\begin{center}
\begin{tabular}{lp{8cm}}
\textbf{Form} & \textbf{When to use} \\
\hline
\texttt{let} & Independent bindings that do not depend on each other \\
\texttt{let*} & Sequential bindings where later ones depend on earlier ones \\
\texttt{letrec} & Mutually recursive local functions \\
\texttt{letrec*} & Mix of sequential values and recursive functions \\
Named \texttt{let} & Tail-recursive loops with named state \\
Internal \texttt{define} & Local helper functions; equivalent to \texttt{letrec*} \\
\texttt{let-values} & Binding multiple return values \\
\end{tabular}
\end{center}

\section*{Summary}

\begin{itemize}
    \item \texttt{let} binds variables in parallel, while \texttt{let*} binds them sequentially.
    \item \texttt{letrec} and \texttt{letrec*} let bindings refer to one another, as for mutually recursive local functions.
    \item Named \texttt{let} is a concise loop, and internal \texttt{define}s behave like \texttt{letrec*}.
    \item \texttt{let} is syntactic sugar for an immediately applied \texttt{lambda}.
    \item \texttt{let-values} binds the results of an expression that returns multiple values.
\end{itemize}

\section*{Exercises}

\begin{exercise}[Scoping Quiz]
Predict the result of each expression below, then check your answers in the REPL. Pay careful attention to which binding form is used:

\begin{lstlisting}
;; Expression A
(let ((x 1))
  (let ((x 2)
        (y x))
    y))

;; Expression B
(let ((x 1))
  (let* ((x 2)
         (y x))
    y))

;; Expression C
(let ((x 10))
  (let ((x (* x 2))
        (y (+ x 3)))
    (list x y)))
\end{lstlisting}

Why do expressions A and B give different results? What value of \code{x} does \code{y} see in each case?
\end{exercise}

\begin{exercise}[Iterative Factorial]
Write a function \code{factorial} using named \code{let} as an iterative loop with an accumulator. Do not use internal \code{define} for a helper---the named \code{let} should be the entire body:

\begin{lstlisting}[style=repl]
> (factorial 0)
1
> (factorial 5)
120
> (factorial 20)
2432902008176640000
\end{lstlisting}

This should be tail-recursive and handle large inputs thanks to Kaappi's bignum support.
\end{exercise}

\begin{exercise}[Euclidean GCD]
Write a function \code{euclidean-gcd} that computes the greatest common divisor of two positive integers using the Euclidean algorithm. Use named \code{let} for the loop:

\begin{lstlisting}[style=repl]
> (euclidean-gcd 48 18)
6
> (euclidean-gcd 100 75)
25
> (euclidean-gcd 17 13)
1
\end{lstlisting}

The algorithm is: if \code{b} is zero, the GCD is \code{a}; otherwise, replace \code{(a, b)} with \code{(b, a mod b)} and repeat. In Python: \texttt{while b: a, b = b, a \% b}. The named \code{let} version maps each ``iteration'' to a tail call.

Kaappi provides a built-in \code{gcd} procedure---the purpose here is to practice named \code{let}.
\end{exercise}

\begin{exercise}[Mutual Recursion with \code{letrec}]
Write a function \code{balanced-parens?} that checks whether a string of parentheses is balanced (e.g., \code{"(()())"} is balanced, \code{"(()"} is not). Use \code{letrec} to define two mutually recursive helpers: \code{expect-open} looks for an opening paren or end of input, and \code{expect-close} tracks the nesting depth:

\begin{lstlisting}[style=repl]
> (balanced-parens? "(()())")
#t
> (balanced-parens? "(()")
#f
> (balanced-parens? "")
#t
> (balanced-parens? ")(")
#f
\end{lstlisting}

Hint: convert the string to a list of characters with \code{string->list}, then walk the list recursively. One approach uses a single recursive function with a depth counter, but the exercise asks for \code{letrec} with two cooperating functions to practice mutual recursion.
\end{exercise}

In the next chapter, we turn to higher-order functions---procedures that take functions as arguments or return them as results.

\chapterend

\chapter{Higher-Order Functions}

Functions that take other functions as arguments or return them as results are called higher-order functions. This chapter covers \texttt{map}, \texttt{filter}, \texttt{fold}, \texttt{for-each}, and \texttt{apply}---the workhorses of functional programming in Scheme.

\chapter{Input and Output}

This chapter covers Scheme's port-based I/O system: reading from and writing to the console, working with files, and understanding the difference between \texttt{display}, \texttt{write}, and \texttt{newline}. You will learn how to build programs that interact with the outside world.


% Part III: Advanced Topics
\part{Advanced Topics}

\chapter{Macros}
\index{macros}\index{syntax-rules}
\label{ch:macros}

Macros are what make Lisp languages uniquely powerful. They let you extend the language itself---defining new syntactic forms that look and behave exactly like built-in ones. If you find yourself writing the same boilerplate pattern repeatedly, a macro can eliminate it.

Scheme's macro system is \emph{hygienic}\index{hygienic macros}: variables introduced by a macro never accidentally capture variables from the surrounding code. This eliminates an entire class of bugs that plague macro systems in C and Common Lisp.

This chapter covers \texttt{define-syntax} and \texttt{syntax-rules}---the pattern-based macro system specified by R7RS---and, at the end, Kaappi's procedural alternative, \texttt{er-macro-transformer}.

\section{What Is a Macro?}

A macro is a compile-time transformation rule\index{macro expansion}. When Kaappi encounters a macro call, it does not evaluate the arguments---instead, it rewrites the entire expression according to the macro's rules, and then evaluates the result.

One boundary is worth stating up front: macros transform \emph{expressions}, the tree of parenthesized forms the reader has already parsed. They cannot change how characters are read in the first place. Some Lisps let you extend the reader itself with new lexical syntax; Kaappi (like standard R7RS) does not---the token syntax of Chapter~\ref{ch:first-steps} is fixed.\index{reader}

Consider \texttt{when}:

\begin{lstlisting}[style=repl]
> (when (> 5 3)
    (display "yes")
    (newline))
yes
\end{lstlisting}

\texttt{when} is not a built-in---it is a macro that expands to an \texttt{if}:

\begin{lstlisting}
(if (> 5 3)
    (begin
      (display "yes")
      (newline)))
\end{lstlisting}

The macro takes the \emph{source code} and transforms it into different source code before evaluation:

\begin{center}
\begin{tikzpicture}[
    box/.style={rectangle, draw, rounded corners, font=\ttfamily\scriptsize,
                text width=5.8cm, align=left, inner sep=4pt},
    arr/.style={->, >=stealth, thick},
    lbl/.style={font=\sffamily\scriptsize, fill=white, inner sep=1pt}
]
\node[box] (src) {(when (> 5 3)\\~~(display "yes")\\~~(newline))};
\node[box, below=1.2cm of src] (exp) {(if (> 5 3)\\~~(begin\\~~~~(display "yes")\\~~~~(newline)))};
\draw[arr] (src) -- (exp) node[lbl, midway] {macro expansion};
\end{tikzpicture}
\end{center}

\begin{note}[Coming from Python]
Python has no equivalent feature. Decorators modify functions but cannot create new syntax. Macros operate at the level of code itself---they generate code from code.
\end{note}

\section{Your First Macro}

Define a macro with \texttt{define-syntax}\index{define-syntax} and \texttt{syntax-rules}:

\begin{lstlisting}[style=repl]
> (define-syntax my-when
    (syntax-rules ()
      ((my-when test body ...)
       (if test (begin body ...)))))
> (my-when (> 5 3)
    (display "hello")
    (newline))
hello
\end{lstlisting}

Let's break this down:

\begin{itemize}
    \item \texttt{define-syntax} names the macro (\texttt{my-when}).
    \item \texttt{syntax-rules ()} introduces the transformation rules. The \texttt{()} is a list of literal keywords (empty here).
    \item \texttt{(my-when test body ...)} is the \emph{pattern}---it matches the macro call.
    \item \texttt{(if test (begin body ...))} is the \emph{template}---what the macro expands to.
    \item \texttt{...} (ellipsis) means ``zero or more of the preceding element.''
\end{itemize}

\section{Seeing the Expansion}
\index{macro expansion!inspecting}\index{expand REPL command@\texttt{,expand} REPL command}

When a macro misbehaves, look at what it expands to. The Kaappi REPL's \texttt{,expand} command shows the expansion of any expression without evaluating it:

\begin{lstlisting}[style=repl]
> ,expand (when #t (display "yes"))
(if #t (begin (display "yes")))
> ,expand (my-when (> 5 3) (display "hello") (newline))
(if (> 5 3) (begin (display "hello") (newline)))
\end{lstlisting}

Make \texttt{,expand} a habit while working through this chapter: after each macro you write, expand a sample call and check that the generated code is what you intended. It turns macro debugging from guesswork into reading. The same view exists outside the REPL: \texttt{kaappi expand file.scm}\index{kaappi expand} prints a whole program after macro expansion, which is handy when a macro misbehaves only in the middle of a real file.

\section{Pattern Matching}

\texttt{syntax-rules} patterns work by structural matching:

\begin{itemize}
    \item A bare identifier (like \texttt{test} or \texttt{body}) matches any single expression and binds it as a \emph{pattern variable}.
    \item An identifier followed by \texttt{...} matches zero or more expressions.
    \item Literal parentheses match literal parentheses in the input.
    \item Identifiers listed in the literals list are matched by name, not as pattern variables.
    \item The underscore \texttt{\_}\index{underscore pattern} is a wildcard: it matches any single form without binding it. Because the macro name in a call always matches trivially, real-world macros conventionally write the name position as \texttt{\_}: \texttt{((\_ test body ...) ...)}. This book spells the name out for readability, but expect \texttt{\_} in code you read elsewhere.
\end{itemize}

\subsection{Multiple Clauses}

A macro can have multiple patterns. The first matching pattern wins:

\begin{lstlisting}[style=repl]
> (define-syntax my-and
    (syntax-rules ()
      ((my-and) #t)
      ((my-and x) x)
      ((my-and x rest ...)
       (if x (my-and rest ...) #f))))
> (my-and)
#t
> (my-and 1 2 3)
3
> (my-and 1 #f 3)
#f
\end{lstlisting}

This is a recursive macro: the third clause expands to another call to \texttt{my-and} with fewer arguments, until it reaches a base case.

If no clause matches a call, expansion fails with a compile-time syntax error at the call site. The \texttt{dotimes} macro later in this chapter shows how to catch a misuse deliberately with \texttt{syntax-error}.

\subsection{Literal Keywords}
\label{sec:literal-keywords}

The first argument to \texttt{syntax-rules} is a list of identifiers to match literally rather than as pattern variables:

\begin{lstlisting}[style=repl]
> (define-syntax my-if
    (syntax-rules (then else)
      ((my-if test then consequent else alternate)
       (cond (test consequent)
             (else alternate)))))
> (my-if (> 3 2) then "yes" else "no")
"yes"
\end{lstlisting}

Here, \texttt{then} and \texttt{else} are literal keywords---they must appear in the call exactly as written, not as variable names.

\section{The Ellipsis (\texttt{...})}
\index{ellipsis}

The ellipsis is the key to writing macros that accept variable numbers of subforms.

In a pattern, \texttt{expr ...} matches zero or more expressions. In a template, \texttt{expr ...} replicates the template for each matched element:

\begin{lstlisting}[style=repl]
> (define-syntax my-list
    (syntax-rules ()
      ((my-list) '())
      ((my-list elem rest ...)
       (cons elem (my-list rest ...)))))
> (my-list 1 2 3)
(1 2 3)
\end{lstlisting}

A more useful example---a macro that returns the value of the first expression while evaluating the rest for side effects:

\begin{lstlisting}[style=repl]
> (define-syntax begin0
    (syntax-rules ()
      ((begin0 first rest ...)
       (let ((result first))
         rest ...
         result))))
> (begin0 (+ 1 2)
    (display "side effect")
    (newline))
side effect
3
\end{lstlisting}

\subsection{Nested Ellipsis}

Ellipsis can be nested for macros that process lists of lists:

\begin{lstlisting}[style=repl]
> (define-syntax my-let
    (syntax-rules ()
      ((my-let ((var init) ...) body ...)
       ((lambda (var ...) body ...) init ...))))
> (my-let ((x 5) (y 10))
    (+ x y))
15
\end{lstlisting}

This \texttt{my-let} macro transforms into a lambda application---exactly how \texttt{let} is defined in the language.

\section{Hygiene}
\label{sec:hygiene}
\index{hygienic macros}

Scheme macros are hygienic: temporary variables introduced by a macro cannot clash with variables in the calling code. Consider:

\begin{lstlisting}[style=repl]
> (define-syntax swap!
    (syntax-rules ()
      ((swap! a b)
       (let ((tmp a))
         (set! a b)
         (set! b tmp)))))
> (let ((tmp 100) (x 1) (y 2))
    (swap! x y)
    (list tmp x y))
(100 2 1)
\end{lstlisting}

The \texttt{tmp} inside the macro and the \texttt{tmp} in the calling code are different bindings---the macro's \texttt{tmp} is automatically renamed to avoid capture\index{variable capture}. In an unhygienic system (like C's \texttt{\#define}), this would be a bug.

There is no magic here, and \texttt{,expand} lets you watch the mechanism work. During expansion, every identifier that comes from the \emph{template} (rather than from the macro call) is renamed to a fresh, globally unique name:

\begin{lstlisting}[style=repl]
> ,expand (swap! x y)
(__hyg_1_let ((__hyg_2_tmp x)) (set! x y) (set! y __hyg_2_tmp))
\end{lstlisting}

Read the expansion alongside the template. The template's \texttt{tmp} became \texttt{\_\_hyg\_2\_tmp}, so it can never collide with the \texttt{tmp} at the call site. The \texttt{x} and \texttt{y} came from the macro \emph{call}, not the template, so they keep their names and refer to the caller's bindings. Even \texttt{let} was renamed---the expander renames every template identifier uniformly, and the compiler recognizes renamed references to standard forms---which means the macro still works if the caller has shadowed \texttt{let} with a local variable. (The numbers in the fresh names are a session-wide counter, so yours will differ.)

Hygiene, then, is just consistent renaming: template identifiers are renamed away from the caller's namespace, call-site identifiers pass through untouched. \index{gensym}

\begin{note}[Why hygiene matters]
In C, the macro \texttt{\#define SWAP(a,b) \{int tmp=a; a=b; b=tmp;\}} breaks if you write \texttt{SWAP(tmp, x)} because the user's \texttt{tmp} and the macro's \texttt{tmp} collide. Scheme's system makes this impossible.
\end{note}

\section{Practical Macros}

\subsection{A \texttt{unless} Macro}

Both \texttt{when} and \texttt{unless} are built-in, but they are defined as macros internally. Here is how \texttt{unless} could be implemented:

\begin{lstlisting}
(define-syntax unless
  (syntax-rules ()
    ((unless test body ...)
     (if (not test) (begin body ...)))))
\end{lstlisting}

\subsection{A \texttt{while} Loop}
\index{while}

Scheme does not have a built-in \texttt{while}, but you can add one:

\begin{lstlisting}[style=repl]
> (define-syntax while
    (syntax-rules ()
      ((while test body ...)
       (let loop ()
         (when test
           body ...
           (loop))))))
> (let ((i 0))
    (while (< i 5)
      (display i) (display " ")
      (set! i (+ i 1))))
0 1 2 3 4
\end{lstlisting}

\subsection{A \texttt{dotimes} Loop}
\index{dotimes}

\begin{lstlisting}[style=repl]
> (define-syntax dotimes
    (syntax-rules ()
      ((dotimes (var count) body ...)
       (let loop ((var 0))
         (when (< var count)
           body ...
           (loop (+ var 1)))))))
> (dotimes (i 5)
    (display (* i i))
    (display " "))
0 1 4 9 16
\end{lstlisting}

A call that does not match the pattern---say \texttt{(dotimes 5 ...)}, forgetting the parentheses around \texttt{(i 5)}---fails with a syntax error at expansion time. You can make such misuse explicit by adding a catch-all clause that expands to \texttt{syntax-error}\index{syntax-error}, an R7RS form that signals an error during expansion:

\begin{lstlisting}
(define-syntax dotimes
  (syntax-rules ()
    ((dotimes (var count) body ...)
     (let loop ((var 0))
       (when (< var count) body ... (loop (+ var 1)))))
    ((dotimes other ...)
     (syntax-error "dotimes: expected (variable count)"))))
\end{lstlisting}

The catch-all clause matches anything the first clause rejected, and the \texttt{syntax-error} in its template stops compilation. The message string also documents, right in the macro definition, what shape of call was expected.

\subsection{An Assertion Macro}

\begin{lstlisting}[style=repl]
> (define-syntax assert
    (syntax-rules ()
      ((assert expr)
       (unless expr
         (error "Assertion failed" 'expr)))))
> (assert (= (+ 1 1) 2))
> (assert (= (+ 1 1) 3))
error[KP3000]: Assertion failed (= (+ 1 1) 3)
\end{lstlisting}

The failing assertion raises an error object whose message is \texttt{"Assertion failed"} and whose irritant is the \emph{source code} of the assertion---\texttt{'expr} in the template quotes it, which is why the report shows the exact expression that failed. A handler can read both parts back; Chapter~\ref{ch:errors} shows how.

\subsection{A Timing Macro}

This macro reports how long an expression takes to evaluate. \texttt{current-second}, from the standard \texttt{(scheme time)} library (Chapter~\ref{ch:stdlib}), returns the current time in seconds:

\begin{lstlisting}
(define-syntax time-it
  (syntax-rules ()
    ((time-it expr)
     (let ((start (current-second)))
       (let ((result expr))
         (let ((elapsed (- (current-second) start)))
           (display "Elapsed: ")
           (display elapsed)
           (display "s")
           (newline)
           result))))))
\end{lstlisting}

\section{Local Macros with \texttt{let-syntax}}
\index{let-syntax}\index{letrec-syntax}

You can define macros with limited scope using \texttt{let-syntax}:

\begin{lstlisting}[style=repl]
> (let-syntax ((double (syntax-rules ()
                         ((double x) (+ x x)))))
    (double 21))
42
\end{lstlisting}

The macro \texttt{double} exists only within the body. For mutually recursive macros, use \texttt{letrec-syntax}.

\begin{warning}[Macro arguments may be evaluated more than once]
A \texttt{syntax-rules} macro substitutes its arguments as \emph{code}, not as computed values. If a pattern variable appears several times in the template, its argument expression runs once for each occurrence. A \texttt{square} macro whose body is \texttt{(* x x)} expands \texttt{(square (read))} into \texttt{(* (read) (read))}, reading its input twice. When an argument may have side effects or is expensive, bind it once with \texttt{let} inside the macro before reusing it.
\end{warning}

\section{Procedural Macros: Explicit Renaming}
\label{sec:er-macros}
\index{er-macro-transformer}\index{explicit renaming}\index{procedural macros}\index{SRFI-211}

Everything so far has been \texttt{syntax-rules}: a pattern on one side, a template on the other, and the expander does the rest. That covers most macros, and it is the only macro system R7RS-small specifies. Sometimes, though, the expansion has to be \emph{computed}---checked for shape, built up one element at a time, or generated from a table---and no template can express that. For those cases Kaappi ships a second macro system: \texttt{er-macro-transformer}, the \emph{explicit renaming} transformer standardized by SRFI-211 and imported as \texttt{(srfi 211 explicit-renaming)}. It is an extension beyond R7RS-small, and it is \emph{not} \texttt{syntax-case}\index{syntax-case}, the R6RS system, which Kaappi does not implement.

An explicit-renaming macro is an ordinary procedure that runs at expansion time. It receives three arguments: \texttt{form}, the whole macro call as plain list data; \texttt{rename}, a procedure that turns a symbol into a hygienic identifier; and \texttt{compare}, a procedure that tests whether two identifiers mean the same thing. Whatever the procedure returns \emph{is} the expansion, so quasiquote (Chapter~\ref{ch:lists}) is the natural tool for building it. Here is \texttt{swap!} again, written by hand:

\begin{lstlisting}[style=repl]
> (import (srfi 211 explicit-renaming))
> (define-syntax swap!
    (er-macro-transformer
      (lambda (form rename compare)
        (let ((a (cadr form))
              (b (caddr form))
              (tmp (rename 'tmp)))
          `(let ((,tmp ,a))
             (set! ,a ,b)
             (set! ,b ,tmp))))))
> (let ((tmp 100) (x 1) (y 2))
    (swap! x y)
    (list tmp x y))
(100 2 1)
\end{lstlisting}

Set this beside the \texttt{syntax-rules} version in Section~\ref{sec:hygiene}. There the template's \texttt{tmp} was renamed automatically; here \texttt{(rename 'tmp)} does the same job explicitly, producing the kind of fresh identifier the \texttt{,expand} output showed. The pieces that came from the call, \texttt{a} and \texttt{b}, are spliced in untouched, exactly as pattern variables are. Renaming the same symbol twice within one expansion yields the same identifier, so a binding site and its references agree. A fully defensive version would also write \texttt{(rename 'let)} and \texttt{(rename 'set!)}, so that the macro survives a caller who has shadowed those names; \texttt{syntax-rules} did that for every template identifier without being asked, and that is the convenience you give up here.

\begin{note}[The transformer runs at expansion time]
The \texttt{lambda} inside \texttt{er-macro-transformer} is evaluated when the \texttt{define-syntax} is, and called once per use of the macro---before any of the surrounding code runs. It cannot see runtime variables, and an \texttt{error} raised inside it is reported as a syntax error at the macro's use site rather than as a runtime exception.
\end{note}

\subsection{\texttt{compare}: Matching Keywords}

\texttt{compare} answers the question a literal in Section~\ref{sec:literal-keywords} answers: do these two identifiers refer to the same binding? Hand it a piece of the input and a renamed keyword:

\begin{lstlisting}[style=repl]
> (define-syntax my-if
    (er-macro-transformer
      (lambda (form rename compare)
        (let ((test (list-ref form 1))
              (kw1 (list-ref form 2))
              (yes (list-ref form 3))
              (kw2 (list-ref form 4))
              (no (list-ref form 5)))
          (unless (and (compare kw1 (rename 'then))
                       (compare kw2 (rename 'else)))
            (error "my-if: expected then/else keywords"))
          `(if ,test ,yes ,no)))))
> (my-if (> 3 2) then "yes" else "no")
"yes"
\end{lstlisting}

The check is binding-aware, exactly like a \texttt{syntax-rules} literal. If the caller has a local variable named \texttt{else} in scope at the use site, \texttt{(compare kw2 (rename 'else))} returns \texttt{\#f}: the caller's \texttt{else} and the macro's \texttt{else} are different things, and the macro refuses the form just as a \texttt{syntax-rules} version with \texttt{(then else)} as its literals would.

\subsection{Validation with Good Errors}

Because the transformer sees the whole form, it can check the shape and refuse bad input at expansion time, with a message that says what was wrong:

\begin{lstlisting}[style=repl]
> (define-syntax strict-when
    (er-macro-transformer
      (lambda (form rename compare)
        (if (= (length form) 3)
            `(if ,(cadr form) ,(caddr form))
            (error "strict-when: one body form, got"
                   (- (length form) 2))))))
> (strict-when (> 3 2) 'ok)
ok
\end{lstlisting}

Save the import and the definition as \texttt{strict.scm}, add a good use whose result you \texttt{display}, and then a bad one: \texttt{(strict-when \#t 1 2 3)}. Running the file prints the good result and then stops at the bad use, naming the file, the line and column of that use, and your own message with its irritant:

\begin{Verbatim}[fontsize=\footnotesize]
$ kaappi strict.scm
ok
strict.scm:13:1: syntax-error[KP2002]: strict-when: one body form, got 3
\end{Verbatim}

A \texttt{syntax-rules} macro handed a form none of its clauses match can only report \texttt{compile error[KP2001]: invalid syntax}. Here the message is whatever you wrote, and it appears once, at the offending use, before anything runs.

\subsection{When to Reach for It}

Reach for \texttt{syntax-rules} first: it is shorter, declarative, and portable to every R7RS Scheme. Reach for \texttt{er-macro-transformer} when the expansion must be computed---validated with a precise message, built by deciding something per element (the kind of recursion a template cannot express), or generated from data your program holds. The cost is symmetrical: you construct the output \texttt{syntax-rules} would have constructed for you, and you take over renaming every identifier the expansion introduces. Portable code can test for the library itself:

\begin{lstlisting}[style=repl]
> (cond-expand
    ((library (srfi 211 explicit-renaming)) 'procedural)
    (else 'none))
procedural
\end{lstlisting}

\section{When to Use Macros}

Macros are powerful but should not be overused. Use a macro when:

\begin{itemize}
    \item You need to control \emph{when} or \emph{whether} arguments are evaluated (like \texttt{when}, \texttt{unless}, \texttt{and}).
    \item You need to introduce new bindings (like \texttt{let}, \texttt{dotimes}).
    \item You want to generate code at compile time to avoid runtime overhead.
    \item You are building a domain-specific language (DSL) with custom syntax.
\end{itemize}

Do \emph{not} use a macro when a function would suffice. Functions are simpler, composable (you can pass them to \texttt{map}), and easier to debug. If your ``macro'' does not need to delay evaluation or introduce bindings, make it a function.

\section*{Summary}

\begin{itemize}
    \item Macros transform code at expansion time, before evaluation, letting you add new syntactic forms.
    \item \texttt{syntax-rules} matches input against patterns and rewrites it using output templates.
    \item The ellipsis (\texttt{...}) matches and expands sequences of forms and can be nested.
    \item Scheme macros are hygienic, so introduced identifiers never accidentally capture the user's variables.
    \item The REPL's \texttt{,expand} command shows what a call expands to---the first tool to reach for when a macro misbehaves.
    \item \texttt{er-macro-transformer} (SRFI-211) writes the expansion as a procedure: \texttt{rename} supplies hygiene and \texttt{compare} matches keywords. Use it when the expansion must be computed or validated; otherwise stay with \texttt{syntax-rules}.
    \item \texttt{let-syntax} defines local macros; reach for macros only when a function cannot do the job.
\end{itemize}

\section*{Exercises}

\begin{exercise}[Swap!]
Write a \code{swap!} macro that swaps the values of two variables using \code{set!}. Verify that hygiene prevents the macro's temporary variable from clashing with user code.

\begin{lstlisting}[style=repl]
> (let ((a 1) (b 2))
    (swap! a b)
    (list a b))
(2 1)
> (let ((tmp 99) (x 10) (y 20))
    (swap! x y)
    (list tmp x y))
(99 20 10)
\end{lstlisting}
\end{exercise}

\begin{exercise}[Assert-Equal]
Write an \code{assert-equal} macro that takes two expressions, evaluates both, and raises an error showing the source code of each expression and their values if they are not \code{equal?}. (Hint: use \code{'expr} in the template to capture source code.)

\begin{lstlisting}[style=repl]
> (assert-equal (+ 1 2) 3)
;; no output --- assertion passed
> (assert-equal (* 2 3) 7)
error[KP3000]: assertion failed: (* 2 3) evaluated to 6, expected 7
\end{lstlisting}
\end{exercise}

\begin{exercise}[Do-Times with Named Counter]
Write a \code{do-n-times} macro similar to the \code{dotimes} macro from this chapter, but the user supplies the counter variable name, the count, \emph{and} an optional step size (defaulting to 1). You will need two \code{syntax-rules} clauses.

\begin{lstlisting}[style=repl]
> (do-n-times (i 5)
    (display i) (display " "))
0 1 2 3 4
> (do-n-times (i 10 3)
    (display i) (display " "))
0 3 6 9
\end{lstlisting}
\end{exercise}

\begin{exercise}[With-Logging]
Write a \code{with-logging} macro that prints an entry message before evaluating the body and an exit message after, then returns the body's result. The macro should accept a label string as its first argument.

\begin{lstlisting}[style=repl]
> (with-logging "compute"
    (+ 1 2))
[enter: compute]
[exit: compute]
3
\end{lstlisting}

As a bonus, make it print the elapsed time in the exit message using \code{current-second}.
\end{exercise}

\begin{exercise}[Counting at Expansion Time]
Write \code{count-args}, an explicit-renaming macro that expands to the number of subforms it was given. The transformer computes the count, so the expansion is a literal integer.

\begin{lstlisting}[style=repl]
> (count-args a b c)
3
> (count-args)
0
\end{lstlisting}

Use \code{,expand} to confirm that \code{(count-args a b c)} expands to plain \code{3}, then explain why the arguments are never evaluated.
\end{exercise}

In the next chapter (Chapter~\ref{ch:libraries}), we cover the library system---how to organize code into modules with clean interfaces.

\chapterend

\chapter{The Library System}

R7RS defines a module system based on \texttt{define-library} and \texttt{import}. This chapter explains how to organize code into libraries, control what is exported, use import modifiers (\texttt{only}, \texttt{except}, \texttt{rename}, \texttt{prefix}), and write \texttt{.sld} files that Kaappi can load.

\chapter{Error Handling}

Programs need to handle errors gracefully. Scheme provides a condition system with \texttt{guard}, \texttt{raise}, and \texttt{with-exception-handler}. This chapter covers raising and catching exceptions, defining error types, and writing robust code that fails predictably.

\chapter{Continuations}

Continuations are Scheme's most powerful---and most unusual---control flow mechanism. With \texttt{call/cc}, you can capture ``the rest of the computation'' as a first-class value and invoke it later. This chapter demystifies continuations with practical examples, covers \texttt{dynamic-wind}, and explains escape continuations (\texttt{call/ec}) for simpler use cases.

\chapter{Records and Data Abstraction}

When lists are not structured enough, records provide named fields with type safety. This chapter covers \texttt{define-record-type}, constructors, predicates, accessors, and how to build clean abstractions over your data.


% Part IV: Beyond R7RS
\part{Beyond R7RS}

\chapter{The Standard Library}

Beyond R7RS-small's 14 libraries, Kaappi provides 72 SRFIs (Scheme Requests for Implementation) covering data structures, string processing, sorting, random numbers, pattern matching, and more. This chapter surveys the most useful SRFIs and shows how to import them.

\chapter{Foreign Function Interface}
\index{FFI}\index{foreign function interface}
\label{ch:ffi}

Sometimes you need to use a library written in C---a graphics toolkit, a database driver, a compression algorithm, or a system API. Kaappi's Foreign Function Interface (FFI) lets you call C functions directly from Scheme without writing any C code yourself. This chapter covers opening shared libraries, binding functions, handling types, and passing Scheme callbacks into C.

\section{The \texttt{(kaappi ffi)} Library}

The FFI is provided by the \texttt{(kaappi ffi)}\index{kaappi ffi} library:

\begin{lstlisting}
(import (kaappi ffi))
\end{lstlisting}

The core procedures are:

\begin{center}
\begin{tabular}{lp{6.5cm}}
\textbf{Procedure} & \textbf{Purpose} \\
\hline
\texttt{ffi-open} & Open a shared library (.dylib / .so) \\
\texttt{ffi-fn} & Bind a C function from the library \\
\texttt{ffi-close} & Close the library handle \\
\texttt{ffi-callback} & Create a C function pointer from a Scheme procedure \\
\texttt{ffi-callback-release} & Free a callback \\
\texttt{ffi-callback?} & Test if a value is an FFI callback \\
\texttt{ffi-bytevector-ptr} & Get the data pointer of a bytevector (for passing buffers to C) \\
\end{tabular}
\end{center}

\section{Opening a Shared Library}
\index{ffi-open}

\begin{lstlisting}[style=repl]
> (import (kaappi ffi))
> (define libm (ffi-open "libm.dylib"))  ; macOS
\end{lstlisting}

On Linux, use the \texttt{.so}\index{shared library} name:

\begin{lstlisting}
(define libm (ffi-open "libm.so.6"))
\end{lstlisting}

The library is loaded into the process and stays open until you call \texttt{ffi-close} or the program exits.

\begin{note}[Library names by platform]
macOS uses \texttt{.dylib} extensions (\texttt{libm.dylib}, \texttt{libsqlite3.dylib}). Linux uses \texttt{.so} with version suffixes (\texttt{libm.so.6}, \texttt{libsqlite3.so.0}). You can also pass an absolute path.
\end{note}

\section{Binding C Functions}
\index{ffi-fn}

\texttt{ffi-fn} looks up a symbol in the library and creates a callable Scheme procedure:

\begin{lstlisting}
(ffi-fn library "function-name" '(param-types ...) 'return-type)
\end{lstlisting}

A function can take at most 16 parameters.

\subsection{Example: Math Functions}

\begin{lstlisting}[style=repl]
> (define c-sqrt (ffi-fn libm "sqrt" '(double) 'double))
> (define c-pow (ffi-fn libm "pow" '(double double) 'double))
> (define c-ceil (ffi-fn libm "ceil" '(double) 'double))
> (c-sqrt 2.0)
1.4142135623730951
> (c-pow 2.0 10.0)
1024.0
> (c-ceil 3.2)
4.0
\end{lstlisting}

The bound function behaves like any other Scheme procedure---you can pass it to \texttt{map}, store it in a variable, or use it in higher-order patterns.

\subsection{Example: String Functions}

\begin{lstlisting}[style=repl]
> (define libc (ffi-open "libc.dylib"))
> (define c-strlen (ffi-fn libc "strlen" '(string) 'size_t))
> (c-strlen "hello")
5
> (c-strlen "")
0
\end{lstlisting}

\section{Supported Types}
\index{FFI types}

The following C types are supported in parameter and return type positions:

\begin{center}
\begin{tabular}{ll}
\textbf{FFI type} & \textbf{C type} \\
\hline
\texttt{int} & \texttt{int} (32-bit signed) \\
\texttt{long} & \texttt{long} (64-bit on 64-bit systems) \\
\texttt{int8} / \texttt{uint8} & \texttt{int8\_t} / \texttt{uint8\_t} \\
\texttt{int16} / \texttt{uint16} & \texttt{int16\_t} / \texttt{uint16\_t} \\
\texttt{int32} / \texttt{uint32} & \texttt{int32\_t} / \texttt{uint32\_t} \\
\texttt{int64} / \texttt{uint64} & \texttt{int64\_t} / \texttt{uint64\_t} \\
\texttt{size\_t} & \texttt{size\_t} \\
\texttt{float} & \texttt{float} (32-bit) \\
\texttt{double} & \texttt{double} (64-bit) \\
\texttt{char} & \texttt{char} \\
\texttt{bool} & \texttt{\_Bool} \\
\texttt{string} & \texttt{const char*} (null-terminated) \\
\texttt{pointer} & \texttt{void*} (opaque pointer) \\
\texttt{void} & \texttt{void} (return type only) \\
\end{tabular}
\end{center}

\subsection{Type Conversions}

\begin{itemize}
    \item \textbf{Strings:} Scheme strings are automatically converted to null-terminated C strings when passed as arguments, and C strings are converted back to Scheme strings when returned.
    \item \textbf{Numbers:} Scheme integers and floats are converted to the appropriate C numeric type. Out-of-range values are rejected with an error naming the FFI function and the expected type.
    \item \textbf{Characters:} The \texttt{char} type takes and returns Scheme characters---a C function declared \texttt{(char) -> char} maps \texttt{\#\textbackslash a} to \texttt{\#\textbackslash A}, not \texttt{97} to \texttt{65}.
    \item \textbf{Pointers:}\index{pointer} Opaque pointer values are represented as Scheme FFI pointer objects. You cannot dereference them directly from Scheme.
    \item \textbf{Void:} Use \texttt{'void} as the return type for C functions that return nothing.
\end{itemize}

\section{A Complete Example: System Information}

Here is a more realistic example that combines several libc functions to query system information:

\begin{lstlisting}
(import (kaappi ffi))

(define libc (ffi-open "libc.dylib"))

;; Bind functions we need
(define c-getenv
  (ffi-fn libc "getenv" '(string) 'string))
(define c-strlen
  (ffi-fn libc "strlen" '(string) 'size_t))
(define c-time
  (ffi-fn libc "time" '(pointer) 'long))
(define c-getpid
  (ffi-fn libc "getpid" '() 'int))
\end{lstlisting}

\begin{lstlisting}[style=repl]
> (c-getenv "HOME")
"/Users/alice"    (*@\emph{(your values will differ)}@*)
> (c-getenv "SHELL")
"/bin/zsh"
> (c-strlen (c-getenv "HOME"))
12
> (c-getpid)
42137
> (c-time 0)
1751234567
\end{lstlisting}

\begin{lstlisting}
;; Build a system report
(define (system-info)
  (list
    (cons "home" (c-getenv "HOME"))
    (cons "shell" (c-getenv "SHELL"))
    (cons "pid" (c-getpid))
    (cons "time" (c-time 0))))

;; Clean up
(ffi-close libc)
\end{lstlisting}

\begin{note}[Why not SQLite?]
The SQLite C API uses pointer-to-pointer out-parameters (\texttt{sqlite3**}) to return database handles, which requires low-level pointer manipulation beyond what the FFI provides directly. In practice, you would use the \texttt{kaappi-sqlite} package (\texttt{thottam install kaappi-sqlite}), which wraps these details in a clean Scheme interface.
\end{note}

\section{Passing Buffers with Bytevectors}
\index{ffi-bytevector-ptr}\index{bytevectors!and FFI}

Many C functions do not return a string---they fill a buffer you hand them. \texttt{ffi-bytevector-ptr} returns the address of a bytevector's data, which you can pass wherever C expects a pointer. Here \texttt{getcwd} writes the current directory into a bytevector as a 0-terminated C string:

\begin{lstlisting}
(import (kaappi ffi))

(define libc (ffi-open "libc.dylib"))
(define c-getcwd
  (ffi-fn libc "getcwd" '(pointer size_t) 'pointer))

;; Extract the C string that ends at the first 0 byte
(define (bytevector-cstring bv)
  (let loop ((i 0))
    (if (zero? (bytevector-u8-ref bv i))
        (utf8->string (bytevector-copy bv 0 i))
        (loop (+ i 1)))))

(define buf (make-bytevector 4096 0))
(define result (c-getcwd (ffi-bytevector-ptr buf) 4096))
(display (bytevector-cstring buf))
\end{lstlisting}

\begin{codeoutput}
/Users/alice/projects
\end{codeoutput}

Keep the bytevector reachable for as long as C holds the pointer. Kaappi's garbage collector does not move objects, so the address stays valid while the bytevector is alive---but if the bytevector is collected, the pointer dangles.

\section{Callbacks: Passing Scheme to C}
\index{ffi-callback}

Some C APIs expect function pointers as arguments (e.g., comparators for \texttt{qsort}, event handlers, iteration callbacks). \texttt{ffi-callback} creates a C-compatible function pointer from a Scheme procedure:

\begin{lstlisting}
(define cb (ffi-callback scheme-proc '(param-types ...) 'return-type))
\end{lstlisting}

Callbacks support only a fixed set of signatures:

\begin{center}
\begin{tabular}{ll}
\textbf{Parameters} & \textbf{Return type} \\
\hline
\texttt{()} & \texttt{void} \\
\texttt{(int)} & \texttt{void} \\
\texttt{(pointer)} & \texttt{void} \\
\texttt{(pointer)} & \texttt{int} \\
\texttt{(int pointer)} & \texttt{int} \\
\texttt{(pointer pointer)} & \texttt{void} \\
\texttt{(pointer pointer)} & \texttt{int} \\
\end{tabular}
\end{center}

Any other combination raises \texttt{unsupported callback signature}. At most 32 callbacks can exist at once, which is another reason to release them promptly.

\subsection{Example: Sorting Bytes with \texttt{qsort}}

Pointer arguments arrive in the Scheme procedure as raw addresses. You cannot dereference them from Scheme---but you can hand them straight back to C. This comparator sorts the bytes of a bytevector with C's \texttt{qsort} by delegating each comparison to \texttt{memcmp}:

\begin{lstlisting}
(define c-qsort
  (ffi-fn libc "qsort"
    '(pointer size_t size_t pointer) 'void))
(define c-memcmp
  (ffi-fn libc "memcmp" '(pointer pointer size_t) 'int))

;; qsort passes addresses of two elements; compare 1 byte each
(define byte-compare
  (ffi-callback
    (lambda (a b) (c-memcmp a b 1))
    '(pointer pointer)
    'int))

(define data (bytevector 3 1 4 1 5 9 2 6))
(c-qsort (ffi-bytevector-ptr data) 8 1 byte-compare)
(ffi-callback-release byte-compare)
(display data)
\end{lstlisting}

\begin{codeoutput}
\#u8(1 1 2 3 4 5 6 9)
\end{codeoutput}

\begin{warning}[Always release callbacks]
Callbacks pin Scheme closures so the garbage collector cannot reclaim them. Always call \texttt{ffi-callback-release}\index{ffi-callback-release} when the callback is no longer needed to prevent memory leaks.
\end{warning}

If the Scheme procedure raises an error while C is calling it, the error is not lost: it is re-raised in your program as soon as the C call that invoked the callback returns, where \texttt{guard} can catch it as usual. The C code itself never sees a Scheme exception---it runs to completion first.

\section{Closing Libraries}
\index{ffi-close}

When you are done with a library, close it to release the handle:

\begin{lstlisting}[style=repl]
> (ffi-close libm)
\end{lstlisting}

After closing, any functions bound from that library become invalid. In practice, most programs open libraries at startup and never close them (the OS reclaims everything on exit).

\section{Error Handling}

FFI calls can crash your program if you get the types wrong---there is no safety net between Scheme and C. Common pitfalls:

\begin{itemize}
    \item \textbf{Wrong types:} Passing an integer where a pointer is expected (or vice versa) causes undefined behavior.
    \item \textbf{Null pointers:} C functions may return NULL; check pointer values before using them.
    \item \textbf{Lifetime issues:} If you pass a Scheme string to C and the C code stores the pointer for later use, the Scheme GC may collect the string. Copy the data on the C side if it needs to persist.
    \item \textbf{Thread safety:} Do not call FFI functions from multiple Kaappi fibers simultaneously unless the C library is thread-safe.
\end{itemize}

\begin{warning}[FFI bypasses safety]
The FFI is inherently unsafe. Incorrect type declarations, buffer overflows, and use-after-free in C code can crash the Kaappi process. Test FFI bindings carefully and prefer ecosystem packages (\texttt{kaappi-sqlite}, \texttt{kaappi-pg}, etc.) when available.
\end{warning}

The FFI error messages you are most likely to encounter, with their fixes, are listed in Appendix~\ref{app:errors}.

\subsection{Debugging FFI Problems}

When FFI calls misbehave, use this checklist:

\begin{enumerate}
    \item \textbf{Verify the library loaded.} If \texttt{ffi-open} returns without error, the library is found. If it raises an error, check the file name and path. On macOS, use \texttt{otool -L} to inspect a binary's dependencies; on Linux, use \texttt{ldd}.
    \item \textbf{Check the symbol name.} C++ name-mangling, static linking, and version suffixes can cause \texttt{ffi-fn} to fail. Use \texttt{nm -gU libfoo.dylib | grep function\_name} to confirm the symbol exists.
    \item \textbf{Match types exactly.} A function declared as returning \texttt{int} but actually returning \texttt{size\_t} may work on some platforms and silently corrupt values on others. Consult the C header file.
    \item \textbf{Test in isolation.} Bind one function at a time in the REPL. Call it with known inputs and compare against expected outputs before integrating into a larger program.
\end{enumerate}

\section{FFI and Sandbox Mode}
\index{sandbox}\index{--sandbox}

Kaappi's \texttt{--sandbox} mode disables the FFI entirely. This is useful for running untrusted code:

\begin{lstlisting}
kaappi --sandbox untrusted.scm
\end{lstlisting}

In sandbox mode, \texttt{(import (kaappi ffi))} itself fails with \texttt{sandbox: cannot load library from file}, and none of the FFI procedures are defined. See Appendix~\ref{app:cli} for the full set of command-line options, including \texttt{--sandbox}.

\section{When to Use the FFI}

\begin{itemize}
    \item \textbf{Use it} when you need a C library that has no Scheme wrapper (niche libraries, system APIs, hardware interfaces).
    \item \textbf{Prefer ecosystem packages} (\texttt{thottam install ...}) when they exist---they handle error checking, memory management, and provide idiomatic Scheme interfaces.
    \item \textbf{Consider writing a library} if you find yourself wrapping more than a few functions. Publish it via thottam so others benefit.
\end{itemize}

\section*{Summary}

\begin{itemize}
    \item The \texttt{(kaappi ffi)} library calls C functions in shared libraries at runtime.
    \item \texttt{ffi-open} loads a shared library and \texttt{ffi-fn} binds a C function to a Scheme procedure given its type signature.
    \item The FFI supports integer, floating-point, string, and pointer types; bytevectors pass buffers via \texttt{ffi-bytevector-ptr}, and Scheme procedures become C function pointers via \texttt{ffi-callback} (fixed signature set, released with \texttt{ffi-callback-release}).
    \item The FFI is inherently unsafe and is disabled by \texttt{--sandbox}; prefer ecosystem packages when they exist.
\end{itemize}

\section*{Exercises}

\begin{exercise}[Binding \texttt{strlen}]
Use the FFI to bind \code{strlen} from the C standard library. Open the appropriate shared library for your platform (\code{libc.dylib} on macOS, \code{libc.so.6} on Linux), bind \code{strlen} with the correct type signature, and test it:

\begin{lstlisting}[style=repl]
> (c-strlen "hello")
5
> (c-strlen "")
0
> (c-strlen "Kaappi Scheme")
13
\end{lstlisting}

What happens if you declare the return type as \code{'int} instead of \code{'size\_t}? Does it still work for short strings? Think about why \code{size\_t} is the correct choice (hint: consider strings longer than $2^{31}$ bytes).
\end{exercise}

\begin{exercise}[Comparing FFI \texttt{pow} with \texttt{expt}]
Bind the C function \code{pow} from \code{libm} using the FFI and write a procedure that compares its results with Scheme's built-in \code{expt}:

\begin{lstlisting}
(define (compare-pow base exp)
  (let ((c-result (c-pow (exact->inexact base)
                         (exact->inexact exp)))
        (s-result (expt base exp)))
    (list c-result s-result (- c-result s-result))))
\end{lstlisting}

Test with several inputs: \code{(compare-pow 2 10)}, \code{(compare-pow 3.0 3.0)}, and \code{(compare-pow 2 0.5)}. When do the results differ? Why might the C version and the Scheme version give slightly different answers? (Hint: think about exact vs.\ inexact arithmetic.)
\end{exercise}

\begin{exercise}[Safe Library Wrapper]
Write a procedure \code{(with-library name proc)} that opens a shared library, passes the handle to \code{proc}, and guarantees that \code{ffi-close} is called even if \code{proc} raises an exception. This is the FFI equivalent of Python's \code{with open(...) as f:} pattern.

\begin{lstlisting}
(define (with-library name proc)
  ;; Your implementation here.
  ;; Use guard or dynamic-wind to ensure cleanup.
  ...)
\end{lstlisting}

\code{guard} (Chapter~\ref{ch:errors}) is sufficient here; \code{dynamic-wind} (Chapter~\ref{ch:continuations}) is the more general tool for paired acquire/release. Test it by binding and calling \code{ceil} from libm:

\begin{lstlisting}[style=repl]
> (with-library "libm.dylib"
    (lambda (lib)
      (let ((c-ceil (ffi-fn lib "ceil"
                      '(double) 'double)))
        (c-ceil 3.2))))
4.0
\end{lstlisting}

Verify that the library is closed after the call returns. What happens if you try to use a function bound from a closed library?
\end{exercise}

In the next chapter, we cover concurrency---fibers, channels, and threads.

\chapterend

\chapter{Concurrency}

Kaappi provides two concurrency models: lightweight green threads (fibers) with channels for message passing, and SRFI-18 OS threads for true parallelism. This chapter covers spawning fibers, communicating via channels, and understanding the trade-offs between cooperative and preemptive concurrency.


% Appendices
\appendix
\part{Appendices}

\chapter{R7RS Quick Reference}

This appendix provides a concise reference to the R7RS-small standard: all syntax forms, standard libraries, and key procedures grouped by category. Use it as a lookup when you need to recall the exact signature or behavior of a built-in.

\chapter{Kaappi CLI Reference}

This appendix documents the \texttt{kaappi} command-line interface: all flags, subcommands (\texttt{compile}, \texttt{run}), environment variables, and REPL comma-commands. It also covers the \texttt{thottam} package manager commands.


% Back matter
\backmatter

\printindex

\end{document}
